% !TEX program = lualatex
% ---------------------------------------------------------------------------------------------
%  Readback note 6: Admissible data: scales, ranks and sizes, and their existence
%
%  Build:  python3 tools/rb.py note <slug>
%
%  One unit of a Lean artifact, rewritten as ordinary mathematics. The prose is hand-written. Lean names,
%  line ranges and counts come from data.tex; quoted Lean statements are copied from the artifact by the build.
%
%  Rules of the text: standard terminology; every symbol is explained where it first appears and listed again
%  in the notation at the end; no letter has two meanings; a proof is complete or says what it takes for granted;
%  Lean names appear only in the LEAN lines and in the concordance.
%
%  The signature rule (the build enforces it): a function, relation or predicate that the artifact defines is
%  written only through its notation macro, which the build makes from the spec (`notation`) with one slot for
%  each explicit argument of the Lean definition. No argument is fixed and dropped, also not "for brevity" and
%  not for one section. Every argument and every hypothesis of a stated declaration stands inside its
%  statement, or in a Setting that the statement cites by number; only implicit and instance arguments (the
%  ambient objects) may stand in the conventions.
% ---------------------------------------------------------------------------------------------
\documentclass[11pt,a4paper]{article}
\ifdefined\pdfvariable\pdfvariable suppressoptionalinfo \numexpr 1+2+4+8\relax\fi
\usepackage[a4paper,hmargin=35mm,top=22mm,bottom=27mm,footskip=12mm,nohead]{geometry}

\usepackage{fontspec}
\directlua{
  luaotfload.add_fallback("gfallback", {"STIXTwoMath-Regular.otf:mode=node;", "DejaVuSans.ttf:mode=node;"})
  luaotfload.add_fallback("gmonofallback", {"DejaVuSansMono.ttf:mode=node;", "STIXTwoMath-Regular.otf:mode=node;"})
}
\setmainfont{NewCM10}[Extension=.otf,UprightFont=*-Book,ItalicFont=*-BookItalic,
  BoldFont=*-Bold,BoldItalicFont=*-BoldItalic,SmallCapsFont=*-Book,SmallCapsFeatures={Letters=SmallCaps},
  RawFeature={fallback=gfallback}]
\setsansfont{NewCMSans10}[Extension=.otf,UprightFont=*-Book,ItalicFont=*-BookOblique,
  BoldFont=*-Bold,BoldItalicFont=*-BoldOblique,RawFeature={fallback=gfallback}]
\setmonofont{JetBrainsMono}[Extension=.otf,UprightFont=*-Regular,ItalicFont=*-Italic,
  BoldFont=*-SemiBold,BoldItalicFont=*-SemiBoldItalic,Scale=MatchLowercase,
  RawFeature={-calt;-liga;fallback=gmonofallback}]
\usepackage{amsmath}
\usepackage{amsthm}
\usepackage[math-style=TeX]{unicode-math}
\setmathfont{NewCMMath-Book.otf}
\emergencystretch=2.5em\relax

\usepackage{xcolor}
\usepackage{tikz}
\usetikzlibrary{calc,arrows.meta,positioning,backgrounds}
\usepackage[protrusion=true,expansion=true]{microtype}
\usepackage{enumitem}
\usepackage{needspace}
\usepackage{booktabs}
\usepackage{array}
\usepackage{longtable}
\usepackage{fancyhdr}
\usepackage{lastpage}
\usepackage{multicol}
\usepackage[hidelinks,pdfusetitle]{hyperref}

\definecolor{ink}{HTML}{0F172A}
\definecolor{inksoft}{HTML}{475569}
\definecolor{inkfaint}{HTML}{94A3B8}
\definecolor{rulegrey}{HTML}{CBD5E1}
\definecolor{arith}{HTML}{C2410C}
\definecolor{est}{HTML}{7C3AED}
\definecolor{ana}{HTML}{0891B2}
\definecolor{ok}{HTML}{047857}
\definecolor{leancol}{HTML}{3730A3}

\title{Admissible data: scales, ranks and sizes, and their existence}
\author{grh-compact}
\date{}

% ------------------------------------------------------------------------------------- macros
\newcommand{\leanfont}{\ttfamily\fontsize{8.6}{10.5}\selectfont\color{leancol}}
\newcommand{\lean}[1]{\csname LN@#1\endcsname}
\newcommand{\fillin}[1]{{\color{arith}\itshape[#1]}}  % a place of the template that is not written yet; the build reports every one that is left
\newcommand{\leanquote}[1]{\csname LQ@#1\endcsname}
% generated by rb.py note; do not edit
\newcommand{\ConcDecls}{98}
\newcommand{\Pin}{72fead10}
\newcommand{\PinText}{that of commit {\leanfont 72fead10}}
\newcommand{\ArtifactName}{grh-compact}
\expandafter\newcommand\csname LN@twd\endcsname{\href{https://github.com/lsernau/grh-lean/blob/72fead102f663d5c7a407964a13dc0cf9d977207/grh-compact/src/New/Data/Different.lean\#L41-L41}{{\leanfont towerDeg}}}
\expandafter\newcommand\csname LN@Nn\endcsname{\href{https://github.com/lsernau/grh-lean/blob/72fead102f663d5c7a407964a13dc0cf9d977207/grh-compact/src/New/Data/Different.lean\#L44-L44}{{\leanfont Parameters.\allowbreak{}N₁}}}
\expandafter\newcommand\csname LN@KP\endcsname{\href{https://github.com/lsernau/grh-lean/blob/72fead102f663d5c7a407964a13dc0cf9d977207/grh-compact/src/New/Data/Different.lean\#L64-L64}{{\leanfont Parameters.\allowbreak{}K}}}
\expandafter\newcommand\csname LN@Mom\endcsname{\href{https://github.com/lsernau/grh-lean/blob/72fead102f663d5c7a407964a13dc0cf9d977207/grh-compact/src/New/Data/Different.lean\#L68-L68}{{\leanfont Parameters.\allowbreak{}Mω}}}
\expandafter\newcommand\csname LN@EP\endcsname{\href{https://github.com/lsernau/grh-lean/blob/72fead102f663d5c7a407964a13dc0cf9d977207/grh-compact/src/New/Data/Different.lean\#L79-L79}{{\leanfont Parameters.\allowbreak{}E}}}
\expandafter\newcommand\csname LN@PlP\endcsname{\href{https://github.com/lsernau/grh-lean/blob/72fead102f663d5c7a407964a13dc0cf9d977207/grh-compact/src/New/Data/Different.lean\#L83-L83}{{\leanfont Parameters.\allowbreak{}Φ}}}
\expandafter\newcommand\csname LN@pin\endcsname{\href{https://github.com/lsernau/grh-lean/blob/72fead102f663d5c7a407964a13dc0cf9d977207/grh-compact/src/New/Data/Different.lean\#L135-L135}{{\leanfont Parameters.\allowbreak{}placeOfInt}}}
\expandafter\newcommand\csname LN@pcl\endcsname{\href{https://github.com/lsernau/grh-lean/blob/72fead102f663d5c7a407964a13dc0cf9d977207/grh-compact/src/New/Data/Different.lean\#L139-L139}{{\leanfont Parameters.\allowbreak{}placeClass}}}
\expandafter\newcommand\csname LN@OmP\endcsname{\href{https://github.com/lsernau/grh-lean/blob/72fead102f663d5c7a407964a13dc0cf9d977207/grh-compact/src/New/Data/Different.lean\#L87-L87}{{\leanfont Parameters.\allowbreak{}bigOmega}}}
\expandafter\newcommand\csname LN@CP\endcsname{\href{https://github.com/lsernau/grh-lean/blob/72fead102f663d5c7a407964a13dc0cf9d977207/grh-compact/src/New/Data/AdmissibleData.lean\#L15-L17}{{\leanfont Parameters.\allowbreak{}c}}}
\expandafter\newcommand\csname LN@Data\endcsname{\href{https://github.com/lsernau/grh-lean/blob/72fead102f663d5c7a407964a13dc0cf9d977207/grh-compact/src/New/Data/AdmissibleData.lean\#L22-L45}{{\leanfont Data}}}
\expandafter\newcommand\csname LQ@Data\endcsname{\leanquoteline{0}{structure Data (P : Parameters) where}\leanquoteline{2}{/-- `logX = log x`, `logB = log b`, `logX₁ = log x₁`: the sizes are chosen on the logarithmic}\leanquoteline{2}{scale. -/}\leanquoteline{2}{(logX logB logX₁ : ℝ) (s m r : ℕ) s\_le : s ≤ 32 * logX}\leanquoteline{2}{thick : Finset \_}\leanquoteline{2}{g : Fin 2 → ℕ}\leanquoteline{2}{piG : 𝓞 P.K}\leanquoteline{2}{scale : 1e12 * P.c ≤ log logB ∧ 100 * log logB ≤ logX ∧ 1e6 * logX ≤ logX₁ ∧ 1e6 * logX₁ ≤ logB ∧}\leanquoteline{4}{6 * log logX₁ ≤ log logB}\leanquoteline{2}{m\_mem : ↑m ∈ Icc (exp (logX / 2)) (2 * exp (logX / 2))}\leanquoteline{2}{r\_mem : ↑r ∈ Icc (exp (5 * logX)) (3 * exp (5 * logX))}\leanquoteline{2}{m\_pow : ∃ e, m = 2 \textasciicircum{} (e + 1) ∧ s = (e + 1) * 22}\leanquoteline{2}{r\_pow : ∃ e, r = 3 \textasciicircum{} (e + 1)}\leanquoteline{2}{thick\_eq\_neg\_one\_mod \{m\} (d : m ∣ P.N₁ s) : m ∣ P.N₀ * 3 \textasciicircum{} (.log 3 s + 5) → ∀ σ ∈ thick, P.placeClass d σ = -1}\leanquoteline{2}{three\_term : let φ := P.placeClass P.two\_pow\_dvd\_N₁; ∀ i j, 0 < i → i < j → j ≤ 324 * s →}\leanquoteline{9}{∀ σ ∈ thick, ∀ τ ∈ thick, ∀ ρ ∈ thick, φ (τ / σ) \textasciicircum{} j = φ (ρ / σ) \textasciicircum{} i → φ (τ / σ) = 1}\leanquoteline{2}{g\_prime i : (g i).Prime}\leanquoteline{2}{g\_modEq\_one i : g i ≡ 1 [MOD P.N₁ s * P.d]}\leanquoteline{2}{dvd\_g\_sub\_one i : P.Mω s \textasciicircum{} 2 ∣ g i - 1}\leanquoteline{2}{g\_injective : g.Injective}\leanquoteline{2}{g\_mem i : ↑(g i) ∈ Icc (exp logB \textasciicircum{} 10) (exp logB \textasciicircum{} 11)}\leanquoteline{2}{norm\_piG : piG * star piG = g 0 * g 1}\leanquoteline{2}{exists\_prime\_window : ∀ U ∈ Icc (100 * logB) (2 \textasciicircum{} s), ∃ p, p.Prime ∧ p ≡ 1 [MOD P.Mω s \textasciicircum{} 2] ∧}\leanquoteline{4}{log p ∈ Icc U (U + √logB)}}
\expandafter\newcommand\csname LN@scx\endcsname{\href{https://github.com/lsernau/grh-lean/blob/72fead102f663d5c7a407964a13dc0cf9d977207/grh-compact/src/New/Data/AdmissibleData.lean\#L52-L52}{{\leanfont Data.\allowbreak{}x}}}
\expandafter\newcommand\csname LN@scxo\endcsname{\href{https://github.com/lsernau/grh-lean/blob/72fead102f663d5c7a407964a13dc0cf9d977207/grh-compact/src/New/Data/AdmissibleData.lean\#L55-L55}{{\leanfont Data.\allowbreak{}x₁}}}
\expandafter\newcommand\csname LN@scb\endcsname{\href{https://github.com/lsernau/grh-lean/blob/72fead102f663d5c7a407964a13dc0cf9d977207/grh-compact/src/New/Data/AdmissibleData.lean\#L58-L58}{{\leanfont Data.\allowbreak{}b}}}
\expandafter\newcommand\csname LN@ED\endcsname{\href{https://github.com/lsernau/grh-lean/blob/72fead102f663d5c7a407964a13dc0cf9d977207/grh-compact/src/New/Data/AdmissibleData.lean\#L62-L62}{{\leanfont Data.\allowbreak{}E}}}
\expandafter\newcommand\csname LN@PlD\endcsname{\href{https://github.com/lsernau/grh-lean/blob/72fead102f663d5c7a407964a13dc0cf9d977207/grh-compact/src/New/Data/AdmissibleData.lean\#L66-L66}{{\leanfont Data.\allowbreak{}Φ}}}
\expandafter\newcommand\csname LN@ND\endcsname{\href{https://github.com/lsernau/grh-lean/blob/72fead102f663d5c7a407964a13dc0cf9d977207/grh-compact/src/New/Data/AdmissibleData.lean\#L69-L69}{{\leanfont Data.\allowbreak{}N₁}}}
\expandafter\newcommand\csname LN@fD\endcsname{\href{https://github.com/lsernau/grh-lean/blob/72fead102f663d5c7a407964a13dc0cf9d977207/grh-compact/src/New/Data/AdmissibleData.lean\#L72-L72}{{\leanfont Data.\allowbreak{}f}}}
\expandafter\newcommand\csname LN@cond\endcsname{\href{https://github.com/lsernau/grh-lean/blob/72fead102f663d5c7a407964a13dc0cf9d977207/grh-compact/src/New/Data/AdmissibleData.lean\#L75-L75}{{\leanfont Data.\allowbreak{}cycloConductor}}}
\expandafter\newcommand\csname LN@rdD\endcsname{\href{https://github.com/lsernau/grh-lean/blob/72fead102f663d5c7a407964a13dc0cf9d977207/grh-compact/src/New/Data/AdmissibleData.lean\#L78-L78}{{\leanfont Data.\allowbreak{}rootDisc}}}
\expandafter\newcommand\csname LN@pinD\endcsname{\href{https://github.com/lsernau/grh-lean/blob/72fead102f663d5c7a407964a13dc0cf9d977207/grh-compact/src/New/Data/AdmissibleData.lean\#L113-L113}{{\leanfont Data.\allowbreak{}placeOfInt}}}
\expandafter\newcommand\csname LN@tops\endcsname{\href{https://github.com/lsernau/grh-lean/blob/72fead102f663d5c7a407964a13dc0cf9d977207/grh-compact/src/New/Data/AdmissibleData.lean\#L127-L127}{{\leanfont Data.\allowbreak{}tops}}}
\expandafter\newcommand\csname LN@gD\endcsname{\href{https://github.com/lsernau/grh-lean/blob/72fead102f663d5c7a407964a13dc0cf9d977207/grh-compact/src/New/Data/AdmissibleData.lean\#L124-L124}{{\leanfont Data.\allowbreak{}gProd}}}
\expandafter\newcommand\csname LN@kD\endcsname{\href{https://github.com/lsernau/grh-lean/blob/72fead102f663d5c7a407964a13dc0cf9d977207/grh-compact/src/New/Data/AdmissibleData.lean\#L82-L82}{{\leanfont Data.\allowbreak{}k}}}
\expandafter\newcommand\csname LN@tD\endcsname{\href{https://github.com/lsernau/grh-lean/blob/72fead102f663d5c7a407964a13dc0cf9d977207/grh-compact/src/New/Data/AdmissibleData.lean\#L85-L85}{{\leanfont Data.\allowbreak{}t}}}
\expandafter\newcommand\csname LN@nD\endcsname{\href{https://github.com/lsernau/grh-lean/blob/72fead102f663d5c7a407964a13dc0cf9d977207/grh-compact/src/New/Data/AdmissibleData.lean\#L116-L116}{{\leanfont Data.\allowbreak{}n}}}
\expandafter\newcommand\csname LN@posD\endcsname{\href{https://github.com/lsernau/grh-lean/blob/72fead102f663d5c7a407964a13dc0cf9d977207/grh-compact/src/New/Data/AdmissibleData.lean\#L119-L119}{{\leanfont Data.\allowbreak{}posRank}}}
\expandafter\newcommand\csname LN@hhD\endcsname{\href{https://github.com/lsernau/grh-lean/blob/72fead102f663d5c7a407964a13dc0cf9d977207/grh-compact/src/New/Data/AdmissibleData.lean\#L92-L92}{{\leanfont Data.\allowbreak{}h₂}}}
\expandafter\newcommand\csname LN@hD\endcsname{\href{https://github.com/lsernau/grh-lean/blob/72fead102f663d5c7a407964a13dc0cf9d977207/grh-compact/src/New/Data/AdmissibleData.lean\#L95-L95}{{\leanfont Data.\allowbreak{}h}}}
\expandafter\newcommand\csname LN@rhg\endcsname{\href{https://github.com/lsernau/grh-lean/blob/72fead102f663d5c7a407964a13dc0cf9d977207/grh-compact/src/New/Data/AdmissibleData.lean\#L98-L98}{{\leanfont Data.\allowbreak{}reducedHalfGap}}}
\expandafter\newcommand\csname LN@rup\endcsname{\href{https://github.com/lsernau/grh-lean/blob/72fead102f663d5c7a407964a13dc0cf9d977207/grh-compact/src/New/Data/AdmissibleData.lean\#L88-L88}{{\leanfont Data.\allowbreak{}roundUp}}}
\expandafter\newcommand\csname LN@aD\endcsname{\href{https://github.com/lsernau/grh-lean/blob/72fead102f663d5c7a407964a13dc0cf9d977207/grh-compact/src/New/Data/AdmissibleData.lean\#L101-L101}{{\leanfont Data.\allowbreak{}a}}}
\expandafter\newcommand\csname LN@bone\endcsname{\href{https://github.com/lsernau/grh-lean/blob/72fead102f663d5c7a407964a13dc0cf9d977207/grh-compact/src/New/Data/AdmissibleData.lean\#L105-L105}{{\leanfont Data.\allowbreak{}b₁}}}
\expandafter\newcommand\csname LN@btwo\endcsname{\href{https://github.com/lsernau/grh-lean/blob/72fead102f663d5c7a407964a13dc0cf9d977207/grh-compact/src/New/Data/AdmissibleData.lean\#L109-L109}{{\leanfont Data.\allowbreak{}b₂}}}
\expandafter\newcommand\csname LN@wt\endcsname{\href{https://github.com/lsernau/grh-lean/blob/72fead102f663d5c7a407964a13dc0cf9d977207/grh-compact/src/New/Data/AdmissibleData.lean\#L131-L131}{{\leanfont Data.\allowbreak{}w}}}
\expandafter\newcommand\csname LN@nw\endcsname{\href{https://github.com/lsernau/grh-lean/blob/72fead102f663d5c7a407964a13dc0cf9d977207/grh-compact/src/New/Data/AdmissibleData.lean\#L134-L134}{{\leanfont Data.\allowbreak{}nominalWeight}}}
\expandafter\newcommand\csname LN@Ask\endcsname{\href{https://github.com/lsernau/grh-lean/blob/72fead102f663d5c7a407964a13dc0cf9d977207/grh-compact/src/New/Data/AdmissibleData.lean\#L138-L139}{{\leanfont Data.\allowbreak{}primeSum}}}
\expandafter\newcommand\csname LN@budg\endcsname{\href{https://github.com/lsernau/grh-lean/blob/72fead102f663d5c7a407964a13dc0cf9d977207/grh-compact/src/New/Data/AdmissibleData.lean\#L143-L143}{{\leanfont Data.\allowbreak{}budget}}}
\expandafter\newcommand\csname LN@jet\endcsname{\href{https://github.com/lsernau/grh-lean/blob/72fead102f663d5c7a407964a13dc0cf9d977207/grh-compact/src/New/Data/AdmissibleData.lean\#L147-L147}{{\leanfont Data.\allowbreak{}jetSize}}}
\expandafter\newcommand\csname LN@UD\endcsname{\href{https://github.com/lsernau/grh-lean/blob/72fead102f663d5c7a407964a13dc0cf9d977207/grh-compact/src/New/Data/AdmissibleData.lean\#L151-L151}{{\leanfont Data.\allowbreak{}dirCost}}}
\expandafter\newcommand\csname LN@hlb\endcsname{\href{https://github.com/lsernau/grh-lean/blob/72fead102f663d5c7a407964a13dc0cf9d977207/grh-compact/src/New/Data/Magnitudes.lean\#L18-L19}{{\leanfont Data.\allowbreak{}HasLogBound}}}
\expandafter\newcommand\csname LN@hlu\endcsname{\href{https://github.com/lsernau/grh-lean/blob/72fead102f663d5c7a407964a13dc0cf9d977207/grh-compact/src/New/Data/Magnitudes.lean\#L38-L38}{{\leanfont Data.\allowbreak{}HasLogUpper}}}
\expandafter\newcommand\csname LN@hll\endcsname{\href{https://github.com/lsernau/grh-lean/blob/72fead102f663d5c7a407964a13dc0cf9d977207/grh-compact/src/New/Data/Magnitudes.lean\#L41-L41}{{\leanfont Data.\allowbreak{}HasLogLower}}}
\expandafter\newcommand\csname LN@pos\endcsname{\href{https://github.com/lsernau/grh-lean/blob/72fead102f663d5c7a407964a13dc0cf9d977207/grh-compact/src/New/Data/Magnitudes.lean\#L43-L43}{{\leanfont Data.\allowbreak{}HasLogBound.\allowbreak{}pos}}}
\expandafter\newcommand\csname LN@upabs\endcsname{\href{https://github.com/lsernau/grh-lean/blob/72fead102f663d5c7a407964a13dc0cf9d977207/grh-compact/src/New/Data/Magnitudes.lean\#L94-L94}{{\leanfont Data.\allowbreak{}hasLogUpper\_of\_abs\_le}}}
\expandafter\newcommand\csname LN@upnonneg\endcsname{\href{https://github.com/lsernau/grh-lean/blob/72fead102f663d5c7a407964a13dc0cf9d977207/grh-compact/src/New/Data/Magnitudes.lean\#L96-L97}{{\leanfont Data.\allowbreak{}hasLogUpper\_of\_le\_of\_nonneg}}}
\expandafter\newcommand\csname LN@upnat\endcsname{\href{https://github.com/lsernau/grh-lean/blob/72fead102f663d5c7a407964a13dc0cf9d977207/grh-compact/src/New/Data/Magnitudes.lean\#L111-L111}{{\leanfont Data.\allowbreak{}hasLogUpper\_of\_le}}}
\expandafter\newcommand\csname LN@lowle\endcsname{\href{https://github.com/lsernau/grh-lean/blob/72fead102f663d5c7a407964a13dc0cf9d977207/grh-compact/src/New/Data/Magnitudes.lean\#L113-L113}{{\leanfont Data.\allowbreak{}hasLogLower\_of\_le}}}
\expandafter\newcommand\csname LN@icc\endcsname{\href{https://github.com/lsernau/grh-lean/blob/72fead102f663d5c7a407964a13dc0cf9d977207/grh-compact/src/New/Data/Magnitudes.lean\#L49-L51}{{\leanfont Data.\allowbreak{}HasLogBound.\allowbreak{}of\_mem\_Icc}}}
\expandafter\newcommand\csname LN@ltBd\endcsname{\href{https://github.com/lsernau/grh-lean/blob/72fead102f663d5c7a407964a13dc0cf9d977207/grh-compact/src/New/Data/Magnitudes.lean\#L117-L117}{{\leanfont Data.\allowbreak{}lt\_of\_hasLogBound}}}
\expandafter\newcommand\csname LN@leBd\endcsname{\href{https://github.com/lsernau/grh-lean/blob/72fead102f663d5c7a407964a13dc0cf9d977207/grh-compact/src/New/Data/Magnitudes.lean\#L121-L121}{{\leanfont Data.\allowbreak{}le\_of\_hasLogBound}}}
\expandafter\newcommand\csname LQ@leBd\endcsname{\leanquoteline{0}{protected lemma le\_of\_hasLogBound : F ≤ G}}
\expandafter\newcommand\csname LN@natlt\endcsname{\href{https://github.com/lsernau/grh-lean/blob/72fead102f663d5c7a407964a13dc0cf9d977207/grh-compact/src/New/Data/Magnitudes.lean\#L123-L123}{{\leanfont Data.\allowbreak{}nat\_lt\_of\_hasLogBound}}}
\expandafter\newcommand\csname LN@natle\endcsname{\href{https://github.com/lsernau/grh-lean/blob/72fead102f663d5c7a407964a13dc0cf9d977207/grh-compact/src/New/Data/Magnitudes.lean\#L125-L125}{{\leanfont Data.\allowbreak{}nat\_le\_of\_hasLogBound}}}
\expandafter\newcommand\csname LN@natcast\endcsname{\href{https://github.com/lsernau/grh-lean/blob/72fead102f663d5c7a407964a13dc0cf9d977207/grh-compact/src/New/Data/Magnitudes.lean\#L53-L54}{{\leanfont Data.\allowbreak{}hasLogBound\_natCast}}}
\expandafter\newcommand\csname LN@bdmul\endcsname{\href{https://github.com/lsernau/grh-lean/blob/72fead102f663d5c7a407964a13dc0cf9d977207/grh-compact/src/New/Data/Magnitudes.lean\#L63-L63}{{\leanfont Data.\allowbreak{}hasLogBound\_mul}}}
\expandafter\newcommand\csname LN@bdadd\endcsname{\href{https://github.com/lsernau/grh-lean/blob/72fead102f663d5c7a407964a13dc0cf9d977207/grh-compact/src/New/Data/Magnitudes.lean\#L65-L66}{{\leanfont Data.\allowbreak{}hasLogBound\_add}}}
\expandafter\newcommand\csname LN@bdf\endcsname{\href{https://github.com/lsernau/grh-lean/blob/72fead102f663d5c7a407964a13dc0cf9d977207/grh-compact/src/New/Data/Magnitudes.lean\#L155-L156}{{\leanfont Data.\allowbreak{}hasLogBound\_f}}}
\expandafter\newcommand\csname LN@cge\endcsname{\href{https://github.com/lsernau/grh-lean/blob/72fead102f663d5c7a407964a13dc0cf9d977207/grh-compact/src/New/Data/Magnitudes.lean\#L21-L22}{{\leanfont Data.\allowbreak{}c\_ge}}}
\expandafter\newcommand\csname LN@round\endcsname{\href{https://github.com/lsernau/grh-lean/blob/72fead102f663d5c7a407964a13dc0cf9d977207/grh-compact/src/New/Data/Magnitudes.lean\#L169-L171}{{\leanfont Data.\allowbreak{}roundUp\_mem\_Icc}}}
\expandafter\newcommand\csname LN@rodd\endcsname{\href{https://github.com/lsernau/grh-lean/blob/72fead102f663d5c7a407964a13dc0cf9d977207/grh-compact/src/New/Data/Magnitudes.lean\#L27-L27}{{\leanfont Data.\allowbreak{}r\_odd}}}
\expandafter\newcommand\csname LN@meven\endcsname{\href{https://github.com/lsernau/grh-lean/blob/72fead102f663d5c7a407964a13dc0cf9d977207/grh-compact/src/New/Data/Magnitudes.lean\#L30-L30}{{\leanfont Data.\allowbreak{}m\_even}}}
\expandafter\newcommand\csname LN@lem\endcsname{\href{https://github.com/lsernau/grh-lean/blob/72fead102f663d5c7a407964a13dc0cf9d977207/grh-compact/src/New/Data/Magnitudes.lean\#L137-L137}{{\leanfont Data.\allowbreak{}le\_m}}}
\expandafter\newcommand\csname LN@radd\endcsname{\href{https://github.com/lsernau/grh-lean/blob/72fead102f663d5c7a407964a13dc0cf9d977207/grh-compact/src/New/Data/Magnitudes.lean\#L139-L139}{{\leanfont Data.\allowbreak{}r\_add}}}
\expandafter\newcommand\csname LN@kler\endcsname{\href{https://github.com/lsernau/grh-lean/blob/72fead102f663d5c7a407964a13dc0cf9d977207/grh-compact/src/New/Data/Magnitudes.lean\#L141-L141}{{\leanfont Data.\allowbreak{}k\_le\_r}}}
\expandafter\newcommand\csname LN@hundred\endcsname{\href{https://github.com/lsernau/grh-lean/blob/72fead102f663d5c7a407964a13dc0cf9d977207/grh-compact/src/New/Data/Magnitudes.lean\#L195-L195}{{\leanfont Data.\allowbreak{}hundred\_mul\_m\_le\_k}}}
\expandafter\newcommand\csname LN@kpos\endcsname{\href{https://github.com/lsernau/grh-lean/blob/72fead102f663d5c7a407964a13dc0cf9d977207/grh-compact/src/New/Data/Magnitudes.lean\#L197-L197}{{\leanfont Data.\allowbreak{}k\_pos}}}
\expandafter\newcommand\csname LN@nmono\endcsname{\href{https://github.com/lsernau/grh-lean/blob/72fead102f663d5c7a407964a13dc0cf9d977207/grh-compact/src/New/Data/Magnitudes.lean\#L199-L199}{{\leanfont Data.\allowbreak{}n\_mono}}}
\expandafter\newcommand\csname LN@keven\endcsname{\href{https://github.com/lsernau/grh-lean/blob/72fead102f663d5c7a407964a13dc0cf9d977207/grh-compact/src/New/Data/Magnitudes.lean\#L201-L201}{{\leanfont Data.\allowbreak{}k\_even}}}
\expandafter\newcommand\csname LN@todd\endcsname{\href{https://github.com/lsernau/grh-lean/blob/72fead102f663d5c7a407964a13dc0cf9d977207/grh-compact/src/New/Data/Magnitudes.lean\#L205-L205}{{\leanfont Data.\allowbreak{}t\_odd}}}
\expandafter\newcommand\csname LN@rsix\endcsname{\href{https://github.com/lsernau/grh-lean/blob/72fead102f663d5c7a407964a13dc0cf9d977207/grh-compact/src/New/Data/Magnitudes.lean\#L207-L207}{{\leanfont Data.\allowbreak{}r\_add\_six\_le\_t}}}
\expandafter\newcommand\csname LN@kth\endcsname{\href{https://github.com/lsernau/grh-lean/blob/72fead102f663d5c7a407964a13dc0cf9d977207/grh-compact/src/New/Data/Magnitudes.lean\#L210-L210}{{\leanfont Data.\allowbreak{}k\_add\_t\_add\_one\_le\_h₂}}}
\expandafter\newcommand\csname LN@clthick\endcsname{\href{https://github.com/lsernau/grh-lean/blob/72fead102f663d5c7a407964a13dc0cf9d977207/grh-compact/src/New/Data/Magnitudes.lean\#L34-L35}{{\leanfont Data.\allowbreak{}placeClass\_of\_thick}}}
\expandafter\newcommand\csname LN@nwicc\endcsname{\href{https://github.com/lsernau/grh-lean/blob/72fead102f663d5c7a407964a13dc0cf9d977207/grh-compact/src/New/Data/Weights.lean\#L14-L16}{{\leanfont Data.\allowbreak{}nominalWeight\_mem\_Icc}}}
\expandafter\newcommand\csname LN@nwnonneg\endcsname{\href{https://github.com/lsernau/grh-lean/blob/72fead102f663d5c7a407964a13dc0cf9d977207/grh-compact/src/New/Data/Weights.lean\#L21-L21}{{\leanfont Data.\allowbreak{}nominalWeight\_nonneg}}}
\expandafter\newcommand\csname LN@nwadd\endcsname{\href{https://github.com/lsernau/grh-lean/blob/72fead102f663d5c7a407964a13dc0cf9d977207/grh-compact/src/New/Data/Weights.lean\#L25-L27}{{\leanfont Data.\allowbreak{}nominalWeight\_add\_le}}}
\expandafter\newcommand\csname LN@nwlew\endcsname{\href{https://github.com/lsernau/grh-lean/blob/72fead102f663d5c7a407964a13dc0cf9d977207/grh-compact/src/New/Data/Weights.lean\#L29-L30}{{\leanfont Data.\allowbreak{}nominalWeight\_le\_w\_sub}}}
\expandafter\newcommand\csname LN@wlenw\endcsname{\href{https://github.com/lsernau/grh-lean/blob/72fead102f663d5c7a407964a13dc0cf9d977207/grh-compact/src/New/Data/Weights.lean\#L36-L37}{{\leanfont Data.\allowbreak{}w\_le\_nominalWeight\_add}}}
\expandafter\newcommand\csname LN@wletwob\endcsname{\href{https://github.com/lsernau/grh-lean/blob/72fead102f663d5c7a407964a13dc0cf9d977207/grh-compact/src/New/Data/Weights.lean\#L40-L41}{{\leanfont Data.\allowbreak{}w\_le\_two\_mul\_b}}}
\expandafter\newcommand\csname LN@logw\endcsname{\href{https://github.com/lsernau/grh-lean/blob/72fead102f663d5c7a407964a13dc0cf9d977207/grh-compact/src/New/Data/Weights.lean\#L51-L54}{{\leanfont Data.\allowbreak{}log\_w\_le}}}
\expandafter\newcommand\csname LN@chequiv\endcsname{\href{https://github.com/lsernau/grh-lean/blob/72fead102f663d5c7a407964a13dc0cf9d977207/grh-compact/src/New/Data/ChosenPlaceLevels.lean\#L18-L19}{{\leanfont Data.\allowbreak{}chosenEquiv}}}
\expandafter\newcommand\csname LN@cardch\endcsname{\href{https://github.com/lsernau/grh-lean/blob/72fead102f663d5c7a407964a13dc0cf9d977207/grh-compact/src/New/Data/ChosenPlaceLevels.lean\#L23-L24}{{\leanfont Data.\allowbreak{}card\_chosenPlace}}}
\expandafter\newcommand\csname LN@prodq\endcsname{\href{https://github.com/lsernau/grh-lean/blob/72fead102f663d5c7a407964a13dc0cf9d977207/grh-compact/src/New/Data/ChosenPlaceLevels.lean\#L38-L39}{{\leanfont Data.\allowbreak{}prod\_q\_chosenPlace}}}
\expandafter\newcommand\csname LN@chstar\endcsname{\href{https://github.com/lsernau/grh-lean/blob/72fead102f663d5c7a407964a13dc0cf9d977207/grh-compact/src/New/Data/ChosenPlaceLevels.lean\#L46-L47}{{\leanfont Data.\allowbreak{}isChosen\_star\_iff}}}
\expandafter\newcommand\csname LN@notstar\endcsname{\href{https://github.com/lsernau/grh-lean/blob/72fead102f663d5c7a407964a13dc0cf9d977207/grh-compact/src/New/Data/ChosenPlaceLevels.lean\#L81-L82}{{\leanfont Data.\allowbreak{}IsChosen.\allowbreak{}not\_star}}}
\expandafter\newcommand\csname LN@rangech\endcsname{\href{https://github.com/lsernau/grh-lean/blob/72fead102f663d5c7a407964a13dc0cf9d977207/grh-compact/src/New/Data/ChosenPlaceLevels.lean\#L59-L60}{{\leanfont Data.\allowbreak{}range\_chosen}}}
\expandafter\newcommand\csname LN@conjapart\endcsname{\href{https://github.com/lsernau/grh-lean/blob/72fead102f663d5c7a407964a13dc0cf9d977207/grh-compact/src/New/Data/ChosenPlaceLevels.lean\#L56-L57}{{\leanfont Data.\allowbreak{}isConjApart\_chosen}}}
\expandafter\newcommand\csname LN@valN\endcsname{\href{https://github.com/lsernau/grh-lean/blob/72fead102f663d5c7a407964a13dc0cf9d977207/grh-compact/src/New/Data/ChosenPlaceLevels.lean\#L26-L28}{{\leanfont Data.\allowbreak{}v\_N\_eq\_exp\_neg\_one}}}
\expandafter\newcommand\csname LN@valpi\endcsname{\href{https://github.com/lsernau/grh-lean/blob/72fead102f663d5c7a407964a13dc0cf9d977207/grh-compact/src/New/Data/ChosenPlaceLevels.lean\#L30-L32}{{\leanfont Data.\allowbreak{}intValuation\_pi}}}
\expandafter\newcommand\csname LN@ordch\endcsname{\href{https://github.com/lsernau/grh-lean/blob/72fead102f663d5c7a407964a13dc0cf9d977207/grh-compact/src/New/Data/ChosenPlaceLevels.lean\#L36-L36}{{\leanfont Data.\allowbreak{}chosenGroupOrder}}}
\expandafter\newcommand\csname LN@chlevel\endcsname{\href{https://github.com/lsernau/grh-lean/blob/72fead102f663d5c7a407964a13dc0cf9d977207/grh-compact/src/New/Data/ChosenPlaceLevels.lean\#L44-L44}{{\leanfont Data.\allowbreak{}chosenLevel}}}
\expandafter\newcommand\csname LN@cardGL\endcsname{\href{https://github.com/lsernau/grh-lean/blob/72fead102f663d5c7a407964a13dc0cf9d977207/grh-compact/src/New/Data/ChosenPlaceLevels.lean\#L51-L54}{{\leanfont Data.\allowbreak{}card\_GL\_dvd\_chosenGroupOrder}}}
\expandafter\newcommand\csname LN@princ\endcsname{\href{https://github.com/lsernau/grh-lean/blob/72fead102f663d5c7a407964a13dc0cf9d977207/grh-compact/src/New/Data/ChosenPlaceLevels.lean\#L62-L63}{{\leanfont Data.\allowbreak{}principalLevel\_le\_chosenLevel}}}
\expandafter\newcommand\csname LN@neat\endcsname{\href{https://github.com/lsernau/grh-lean/blob/72fead102f663d5c7a407964a13dc0cf9d977207/grh-compact/src/New/Data/ChosenPlaceLevels.lean\#L67-L70}{{\leanfont Data.\allowbreak{}classAut\_eq\_bot\_of\_le\_chosenLevel}}}
\expandafter\newcommand\csname LN@mink\endcsname{\href{https://github.com/lsernau/grh-lean/blob/72fead102f663d5c7a407964a13dc0cf9d977207/grh-compact/src/New/Data/ChosenPlaceLevels.lean\#L75-L79}{{\leanfont Data.\allowbreak{}card\_classAut\_dvd\_chosenGroupOrder}}}
\expandafter\newcommand\csname LN@omega\endcsname{\href{https://github.com/lsernau/grh-lean/blob/72fead102f663d5c7a407964a13dc0cf9d977207/grh-compact/src/New/Data/ChosenPlaceLevels.lean\#L86-L90}{{\leanfont Data.\allowbreak{}omega\_zpow\_ne\_one}}}
\expandafter\newcommand\csname LN@behrend\endcsname{\href{https://github.com/lsernau/grh-lean/blob/72fead102f663d5c7a407964a13dc0cf9d977207/grh-compact/src/New/Data/DataExist.lean\#L20-L36}{{\leanfont exists\_three\_term\_free}}}
\expandafter\newcommand\csname LN@pigeon\endcsname{\href{https://github.com/lsernau/grh-lean/blob/72fead102f663d5c7a407964a13dc0cf9d977207/grh-compact/src/New/Data/DataExist.lean\#L41-L45}{{\leanfont exists\_mul\_principal}}}
\expandafter\newcommand\csname LN@exdata\endcsname{\href{https://github.com/lsernau/grh-lean/blob/72fead102f663d5c7a407964a13dc0cf9d977207/grh-compact/src/New/Data/DataExist.lean\#L49-L108}{{\leanfont exists\_data}}}
\expandafter\newcommand\csname LQ@exdata\endcsname{\leanquoteline{0}{lemma exists\_data P \{ε : ℝ\} (e0 : 0 < ε / 6) (e1 : ε < 1 / 100) : ∀ᶠ x in atTop, ∃ D : Data P,}\leanquoteline{4}{D.x = x ∧ log D.x₁ = x \textasciicircum{} (ε / 6) ∧ (1e6 * log x) \textasciicircum{} 2 ≤ log D.x₁ ∧}\leanquoteline{4}{log (log D.b) = ε * log x ∧ towerDeg D.s / log D.b ≤ \#D.thick}}
\newcommand{\ConcordanceRows}{\crow{twd}{\href{https://github.com/lsernau/grh-lean/blob/72fead102f663d5c7a407964a13dc0cf9d977207/grh-compact/src/New/Data/Different.lean\#L41-L41}{{\leanfont towerDeg}}}{Different.lean}{41--41}
\crow{Nn}{\href{https://github.com/lsernau/grh-lean/blob/72fead102f663d5c7a407964a13dc0cf9d977207/grh-compact/src/New/Data/Different.lean\#L44-L44}{{\leanfont Parameters.\allowbreak{}N₁}}}{Different.lean}{44--44}
\crow{KP}{\href{https://github.com/lsernau/grh-lean/blob/72fead102f663d5c7a407964a13dc0cf9d977207/grh-compact/src/New/Data/Different.lean\#L64-L64}{{\leanfont Parameters.\allowbreak{}K}}}{Different.lean}{64--64}
\crow{Mom}{\href{https://github.com/lsernau/grh-lean/blob/72fead102f663d5c7a407964a13dc0cf9d977207/grh-compact/src/New/Data/Different.lean\#L68-L68}{{\leanfont Parameters.\allowbreak{}Mω}}}{Different.lean}{68--68}
\crow{EP}{\href{https://github.com/lsernau/grh-lean/blob/72fead102f663d5c7a407964a13dc0cf9d977207/grh-compact/src/New/Data/Different.lean\#L79-L79}{{\leanfont Parameters.\allowbreak{}E}}}{Different.lean}{79--79}
\crow{PlP}{\href{https://github.com/lsernau/grh-lean/blob/72fead102f663d5c7a407964a13dc0cf9d977207/grh-compact/src/New/Data/Different.lean\#L83-L83}{{\leanfont Parameters.\allowbreak{}Φ}}}{Different.lean}{83--83}
\crow{pin}{\href{https://github.com/lsernau/grh-lean/blob/72fead102f663d5c7a407964a13dc0cf9d977207/grh-compact/src/New/Data/Different.lean\#L135-L135}{{\leanfont Parameters.\allowbreak{}placeOfInt}}}{Different.lean}{135--135}
\crow{pcl}{\href{https://github.com/lsernau/grh-lean/blob/72fead102f663d5c7a407964a13dc0cf9d977207/grh-compact/src/New/Data/Different.lean\#L139-L139}{{\leanfont Parameters.\allowbreak{}placeClass}}}{Different.lean}{139--139}
\crow{OmP}{\href{https://github.com/lsernau/grh-lean/blob/72fead102f663d5c7a407964a13dc0cf9d977207/grh-compact/src/New/Data/Different.lean\#L87-L87}{{\leanfont Parameters.\allowbreak{}bigOmega}}}{Different.lean}{87--87}
\crow{CP}{\href{https://github.com/lsernau/grh-lean/blob/72fead102f663d5c7a407964a13dc0cf9d977207/grh-compact/src/New/Data/AdmissibleData.lean\#L15-L17}{{\leanfont Parameters.\allowbreak{}c}}}{AdmissibleData.lean}{15--17}
\crow{Data}{\href{https://github.com/lsernau/grh-lean/blob/72fead102f663d5c7a407964a13dc0cf9d977207/grh-compact/src/New/Data/AdmissibleData.lean\#L22-L45}{{\leanfont Data}}}{AdmissibleData.lean}{22--45}
\crow{scx}{\href{https://github.com/lsernau/grh-lean/blob/72fead102f663d5c7a407964a13dc0cf9d977207/grh-compact/src/New/Data/AdmissibleData.lean\#L52-L52}{{\leanfont Data.\allowbreak{}x}}}{AdmissibleData.lean}{52--52}
\crow{scxo}{\href{https://github.com/lsernau/grh-lean/blob/72fead102f663d5c7a407964a13dc0cf9d977207/grh-compact/src/New/Data/AdmissibleData.lean\#L55-L55}{{\leanfont Data.\allowbreak{}x₁}}}{AdmissibleData.lean}{55--55}
\crow{scb}{\href{https://github.com/lsernau/grh-lean/blob/72fead102f663d5c7a407964a13dc0cf9d977207/grh-compact/src/New/Data/AdmissibleData.lean\#L58-L58}{{\leanfont Data.\allowbreak{}b}}}{AdmissibleData.lean}{58--58}
\crow{ED}{\href{https://github.com/lsernau/grh-lean/blob/72fead102f663d5c7a407964a13dc0cf9d977207/grh-compact/src/New/Data/AdmissibleData.lean\#L62-L62}{{\leanfont Data.\allowbreak{}E}}}{AdmissibleData.lean}{62--62}
\crow{PlD}{\href{https://github.com/lsernau/grh-lean/blob/72fead102f663d5c7a407964a13dc0cf9d977207/grh-compact/src/New/Data/AdmissibleData.lean\#L66-L66}{{\leanfont Data.\allowbreak{}Φ}}}{AdmissibleData.lean}{66--66}
\crow{ND}{\href{https://github.com/lsernau/grh-lean/blob/72fead102f663d5c7a407964a13dc0cf9d977207/grh-compact/src/New/Data/AdmissibleData.lean\#L69-L69}{{\leanfont Data.\allowbreak{}N₁}}}{AdmissibleData.lean}{69--69}
\crow{fD}{\href{https://github.com/lsernau/grh-lean/blob/72fead102f663d5c7a407964a13dc0cf9d977207/grh-compact/src/New/Data/AdmissibleData.lean\#L72-L72}{{\leanfont Data.\allowbreak{}f}}}{AdmissibleData.lean}{72--72}
\crow{cond}{\href{https://github.com/lsernau/grh-lean/blob/72fead102f663d5c7a407964a13dc0cf9d977207/grh-compact/src/New/Data/AdmissibleData.lean\#L75-L75}{{\leanfont Data.\allowbreak{}cycloConductor}}}{AdmissibleData.lean}{75--75}
\crow{rdD}{\href{https://github.com/lsernau/grh-lean/blob/72fead102f663d5c7a407964a13dc0cf9d977207/grh-compact/src/New/Data/AdmissibleData.lean\#L78-L78}{{\leanfont Data.\allowbreak{}rootDisc}}}{AdmissibleData.lean}{78--78}
\crow{pinD}{\href{https://github.com/lsernau/grh-lean/blob/72fead102f663d5c7a407964a13dc0cf9d977207/grh-compact/src/New/Data/AdmissibleData.lean\#L113-L113}{{\leanfont Data.\allowbreak{}placeOfInt}}}{AdmissibleData.lean}{113--113}
\crow{tops}{\href{https://github.com/lsernau/grh-lean/blob/72fead102f663d5c7a407964a13dc0cf9d977207/grh-compact/src/New/Data/AdmissibleData.lean\#L127-L127}{{\leanfont Data.\allowbreak{}tops}}}{AdmissibleData.lean}{127--127}
\crow{gD}{\href{https://github.com/lsernau/grh-lean/blob/72fead102f663d5c7a407964a13dc0cf9d977207/grh-compact/src/New/Data/AdmissibleData.lean\#L124-L124}{{\leanfont Data.\allowbreak{}gProd}}}{AdmissibleData.lean}{124--124}
\crow{kD}{\href{https://github.com/lsernau/grh-lean/blob/72fead102f663d5c7a407964a13dc0cf9d977207/grh-compact/src/New/Data/AdmissibleData.lean\#L82-L82}{{\leanfont Data.\allowbreak{}k}}}{AdmissibleData.lean}{82--82}
\crow{tD}{\href{https://github.com/lsernau/grh-lean/blob/72fead102f663d5c7a407964a13dc0cf9d977207/grh-compact/src/New/Data/AdmissibleData.lean\#L85-L85}{{\leanfont Data.\allowbreak{}t}}}{AdmissibleData.lean}{85--85}
\crow{nD}{\href{https://github.com/lsernau/grh-lean/blob/72fead102f663d5c7a407964a13dc0cf9d977207/grh-compact/src/New/Data/AdmissibleData.lean\#L116-L116}{{\leanfont Data.\allowbreak{}n}}}{AdmissibleData.lean}{116--116}
\crow{posD}{\href{https://github.com/lsernau/grh-lean/blob/72fead102f663d5c7a407964a13dc0cf9d977207/grh-compact/src/New/Data/AdmissibleData.lean\#L119-L119}{{\leanfont Data.\allowbreak{}posRank}}}{AdmissibleData.lean}{119--119}
\crow{hhD}{\href{https://github.com/lsernau/grh-lean/blob/72fead102f663d5c7a407964a13dc0cf9d977207/grh-compact/src/New/Data/AdmissibleData.lean\#L92-L92}{{\leanfont Data.\allowbreak{}h₂}}}{AdmissibleData.lean}{92--92}
\crow{hD}{\href{https://github.com/lsernau/grh-lean/blob/72fead102f663d5c7a407964a13dc0cf9d977207/grh-compact/src/New/Data/AdmissibleData.lean\#L95-L95}{{\leanfont Data.\allowbreak{}h}}}{AdmissibleData.lean}{95--95}
\crow{rhg}{\href{https://github.com/lsernau/grh-lean/blob/72fead102f663d5c7a407964a13dc0cf9d977207/grh-compact/src/New/Data/AdmissibleData.lean\#L98-L98}{{\leanfont Data.\allowbreak{}reducedHalfGap}}}{AdmissibleData.lean}{98--98}
\crow{rup}{\href{https://github.com/lsernau/grh-lean/blob/72fead102f663d5c7a407964a13dc0cf9d977207/grh-compact/src/New/Data/AdmissibleData.lean\#L88-L88}{{\leanfont Data.\allowbreak{}roundUp}}}{AdmissibleData.lean}{88--88}
\crow{aD}{\href{https://github.com/lsernau/grh-lean/blob/72fead102f663d5c7a407964a13dc0cf9d977207/grh-compact/src/New/Data/AdmissibleData.lean\#L101-L101}{{\leanfont Data.\allowbreak{}a}}}{AdmissibleData.lean}{101--101}
\crow{bone}{\href{https://github.com/lsernau/grh-lean/blob/72fead102f663d5c7a407964a13dc0cf9d977207/grh-compact/src/New/Data/AdmissibleData.lean\#L105-L105}{{\leanfont Data.\allowbreak{}b₁}}}{AdmissibleData.lean}{105--105}
\crow{btwo}{\href{https://github.com/lsernau/grh-lean/blob/72fead102f663d5c7a407964a13dc0cf9d977207/grh-compact/src/New/Data/AdmissibleData.lean\#L109-L109}{{\leanfont Data.\allowbreak{}b₂}}}{AdmissibleData.lean}{109--109}
\crow{wt}{\href{https://github.com/lsernau/grh-lean/blob/72fead102f663d5c7a407964a13dc0cf9d977207/grh-compact/src/New/Data/AdmissibleData.lean\#L131-L131}{{\leanfont Data.\allowbreak{}w}}}{AdmissibleData.lean}{131--131}
\crow{nw}{\href{https://github.com/lsernau/grh-lean/blob/72fead102f663d5c7a407964a13dc0cf9d977207/grh-compact/src/New/Data/AdmissibleData.lean\#L134-L134}{{\leanfont Data.\allowbreak{}nominalWeight}}}{AdmissibleData.lean}{134--134}
\crow{Ask}{\href{https://github.com/lsernau/grh-lean/blob/72fead102f663d5c7a407964a13dc0cf9d977207/grh-compact/src/New/Data/AdmissibleData.lean\#L138-L139}{{\leanfont Data.\allowbreak{}primeSum}}}{AdmissibleData.lean}{138--139}
\crow{budg}{\href{https://github.com/lsernau/grh-lean/blob/72fead102f663d5c7a407964a13dc0cf9d977207/grh-compact/src/New/Data/AdmissibleData.lean\#L143-L143}{{\leanfont Data.\allowbreak{}budget}}}{AdmissibleData.lean}{143--143}
\crow{jet}{\href{https://github.com/lsernau/grh-lean/blob/72fead102f663d5c7a407964a13dc0cf9d977207/grh-compact/src/New/Data/AdmissibleData.lean\#L147-L147}{{\leanfont Data.\allowbreak{}jetSize}}}{AdmissibleData.lean}{147--147}
\crow{UD}{\href{https://github.com/lsernau/grh-lean/blob/72fead102f663d5c7a407964a13dc0cf9d977207/grh-compact/src/New/Data/AdmissibleData.lean\#L151-L151}{{\leanfont Data.\allowbreak{}dirCost}}}{AdmissibleData.lean}{151--151}
\crow{hlb}{\href{https://github.com/lsernau/grh-lean/blob/72fead102f663d5c7a407964a13dc0cf9d977207/grh-compact/src/New/Data/Magnitudes.lean\#L18-L19}{{\leanfont Data.\allowbreak{}HasLogBound}}}{Magnitudes.lean}{18--19}
\crow{hlu}{\href{https://github.com/lsernau/grh-lean/blob/72fead102f663d5c7a407964a13dc0cf9d977207/grh-compact/src/New/Data/Magnitudes.lean\#L38-L38}{{\leanfont Data.\allowbreak{}HasLogUpper}}}{Magnitudes.lean}{38--38}
\crow{hll}{\href{https://github.com/lsernau/grh-lean/blob/72fead102f663d5c7a407964a13dc0cf9d977207/grh-compact/src/New/Data/Magnitudes.lean\#L41-L41}{{\leanfont Data.\allowbreak{}HasLogLower}}}{Magnitudes.lean}{41--41}
\crow{pos}{\href{https://github.com/lsernau/grh-lean/blob/72fead102f663d5c7a407964a13dc0cf9d977207/grh-compact/src/New/Data/Magnitudes.lean\#L43-L43}{{\leanfont Data.\allowbreak{}HasLogBound.\allowbreak{}pos}}}{Magnitudes.lean}{43--43}
\crow{upabs}{\href{https://github.com/lsernau/grh-lean/blob/72fead102f663d5c7a407964a13dc0cf9d977207/grh-compact/src/New/Data/Magnitudes.lean\#L94-L94}{{\leanfont Data.\allowbreak{}hasLogUpper\_of\_abs\_le}}}{Magnitudes.lean}{94--94}
\crow{upnonneg}{\href{https://github.com/lsernau/grh-lean/blob/72fead102f663d5c7a407964a13dc0cf9d977207/grh-compact/src/New/Data/Magnitudes.lean\#L96-L97}{{\leanfont Data.\allowbreak{}hasLogUpper\_of\_le\_of\_nonneg}}}{Magnitudes.lean}{96--97}
\crow{upnat}{\href{https://github.com/lsernau/grh-lean/blob/72fead102f663d5c7a407964a13dc0cf9d977207/grh-compact/src/New/Data/Magnitudes.lean\#L111-L111}{{\leanfont Data.\allowbreak{}hasLogUpper\_of\_le}}}{Magnitudes.lean}{111--111}
\crow{lowle}{\href{https://github.com/lsernau/grh-lean/blob/72fead102f663d5c7a407964a13dc0cf9d977207/grh-compact/src/New/Data/Magnitudes.lean\#L113-L113}{{\leanfont Data.\allowbreak{}hasLogLower\_of\_le}}}{Magnitudes.lean}{113--113}
\crow{icc}{\href{https://github.com/lsernau/grh-lean/blob/72fead102f663d5c7a407964a13dc0cf9d977207/grh-compact/src/New/Data/Magnitudes.lean\#L49-L51}{{\leanfont Data.\allowbreak{}HasLogBound.\allowbreak{}of\_mem\_Icc}}}{Magnitudes.lean}{49--51}
\crow{ltBd}{\href{https://github.com/lsernau/grh-lean/blob/72fead102f663d5c7a407964a13dc0cf9d977207/grh-compact/src/New/Data/Magnitudes.lean\#L117-L117}{{\leanfont Data.\allowbreak{}lt\_of\_hasLogBound}}}{Magnitudes.lean}{117--117}
\crow{leBd}{\href{https://github.com/lsernau/grh-lean/blob/72fead102f663d5c7a407964a13dc0cf9d977207/grh-compact/src/New/Data/Magnitudes.lean\#L121-L121}{{\leanfont Data.\allowbreak{}le\_of\_hasLogBound}}}{Magnitudes.lean}{121--121}
\crow{natlt}{\href{https://github.com/lsernau/grh-lean/blob/72fead102f663d5c7a407964a13dc0cf9d977207/grh-compact/src/New/Data/Magnitudes.lean\#L123-L123}{{\leanfont Data.\allowbreak{}nat\_lt\_of\_hasLogBound}}}{Magnitudes.lean}{123--123}
\crow{natle}{\href{https://github.com/lsernau/grh-lean/blob/72fead102f663d5c7a407964a13dc0cf9d977207/grh-compact/src/New/Data/Magnitudes.lean\#L125-L125}{{\leanfont Data.\allowbreak{}nat\_le\_of\_hasLogBound}}}{Magnitudes.lean}{125--125}
\crow{natcast}{\href{https://github.com/lsernau/grh-lean/blob/72fead102f663d5c7a407964a13dc0cf9d977207/grh-compact/src/New/Data/Magnitudes.lean\#L53-L54}{{\leanfont Data.\allowbreak{}hasLogBound\_natCast}}}{Magnitudes.lean}{53--54}
\crow{bdmul}{\href{https://github.com/lsernau/grh-lean/blob/72fead102f663d5c7a407964a13dc0cf9d977207/grh-compact/src/New/Data/Magnitudes.lean\#L63-L63}{{\leanfont Data.\allowbreak{}hasLogBound\_mul}}}{Magnitudes.lean}{63--63}
\crow{bdadd}{\href{https://github.com/lsernau/grh-lean/blob/72fead102f663d5c7a407964a13dc0cf9d977207/grh-compact/src/New/Data/Magnitudes.lean\#L65-L66}{{\leanfont Data.\allowbreak{}hasLogBound\_add}}}{Magnitudes.lean}{65--66}
\crow{bdf}{\href{https://github.com/lsernau/grh-lean/blob/72fead102f663d5c7a407964a13dc0cf9d977207/grh-compact/src/New/Data/Magnitudes.lean\#L155-L156}{{\leanfont Data.\allowbreak{}hasLogBound\_f}}}{Magnitudes.lean}{155--156}
\crow{cge}{\href{https://github.com/lsernau/grh-lean/blob/72fead102f663d5c7a407964a13dc0cf9d977207/grh-compact/src/New/Data/Magnitudes.lean\#L21-L22}{{\leanfont Data.\allowbreak{}c\_ge}}}{Magnitudes.lean}{21--22}
\crow{round}{\href{https://github.com/lsernau/grh-lean/blob/72fead102f663d5c7a407964a13dc0cf9d977207/grh-compact/src/New/Data/Magnitudes.lean\#L169-L171}{{\leanfont Data.\allowbreak{}roundUp\_mem\_Icc}}}{Magnitudes.lean}{169--171}
\crow{rodd}{\href{https://github.com/lsernau/grh-lean/blob/72fead102f663d5c7a407964a13dc0cf9d977207/grh-compact/src/New/Data/Magnitudes.lean\#L27-L27}{{\leanfont Data.\allowbreak{}r\_odd}}}{Magnitudes.lean}{27--27}
\crow{meven}{\href{https://github.com/lsernau/grh-lean/blob/72fead102f663d5c7a407964a13dc0cf9d977207/grh-compact/src/New/Data/Magnitudes.lean\#L30-L30}{{\leanfont Data.\allowbreak{}m\_even}}}{Magnitudes.lean}{30--30}
\crow{lem}{\href{https://github.com/lsernau/grh-lean/blob/72fead102f663d5c7a407964a13dc0cf9d977207/grh-compact/src/New/Data/Magnitudes.lean\#L137-L137}{{\leanfont Data.\allowbreak{}le\_m}}}{Magnitudes.lean}{137--137}
\crow{radd}{\href{https://github.com/lsernau/grh-lean/blob/72fead102f663d5c7a407964a13dc0cf9d977207/grh-compact/src/New/Data/Magnitudes.lean\#L139-L139}{{\leanfont Data.\allowbreak{}r\_add}}}{Magnitudes.lean}{139--139}
\crow{kler}{\href{https://github.com/lsernau/grh-lean/blob/72fead102f663d5c7a407964a13dc0cf9d977207/grh-compact/src/New/Data/Magnitudes.lean\#L141-L141}{{\leanfont Data.\allowbreak{}k\_le\_r}}}{Magnitudes.lean}{141--141}
\crow{hundred}{\href{https://github.com/lsernau/grh-lean/blob/72fead102f663d5c7a407964a13dc0cf9d977207/grh-compact/src/New/Data/Magnitudes.lean\#L195-L195}{{\leanfont Data.\allowbreak{}hundred\_mul\_m\_le\_k}}}{Magnitudes.lean}{195--195}
\crow{kpos}{\href{https://github.com/lsernau/grh-lean/blob/72fead102f663d5c7a407964a13dc0cf9d977207/grh-compact/src/New/Data/Magnitudes.lean\#L197-L197}{{\leanfont Data.\allowbreak{}k\_pos}}}{Magnitudes.lean}{197--197}
\crow{nmono}{\href{https://github.com/lsernau/grh-lean/blob/72fead102f663d5c7a407964a13dc0cf9d977207/grh-compact/src/New/Data/Magnitudes.lean\#L199-L199}{{\leanfont Data.\allowbreak{}n\_mono}}}{Magnitudes.lean}{199--199}
\crow{keven}{\href{https://github.com/lsernau/grh-lean/blob/72fead102f663d5c7a407964a13dc0cf9d977207/grh-compact/src/New/Data/Magnitudes.lean\#L201-L201}{{\leanfont Data.\allowbreak{}k\_even}}}{Magnitudes.lean}{201--201}
\crow{todd}{\href{https://github.com/lsernau/grh-lean/blob/72fead102f663d5c7a407964a13dc0cf9d977207/grh-compact/src/New/Data/Magnitudes.lean\#L205-L205}{{\leanfont Data.\allowbreak{}t\_odd}}}{Magnitudes.lean}{205--205}
\crow{rsix}{\href{https://github.com/lsernau/grh-lean/blob/72fead102f663d5c7a407964a13dc0cf9d977207/grh-compact/src/New/Data/Magnitudes.lean\#L207-L207}{{\leanfont Data.\allowbreak{}r\_add\_six\_le\_t}}}{Magnitudes.lean}{207--207}
\crow{kth}{\href{https://github.com/lsernau/grh-lean/blob/72fead102f663d5c7a407964a13dc0cf9d977207/grh-compact/src/New/Data/Magnitudes.lean\#L210-L210}{{\leanfont Data.\allowbreak{}k\_add\_t\_add\_one\_le\_h₂}}}{Magnitudes.lean}{210--210}
\crow{clthick}{\href{https://github.com/lsernau/grh-lean/blob/72fead102f663d5c7a407964a13dc0cf9d977207/grh-compact/src/New/Data/Magnitudes.lean\#L34-L35}{{\leanfont Data.\allowbreak{}placeClass\_of\_thick}}}{Magnitudes.lean}{34--35}
\crow{nwicc}{\href{https://github.com/lsernau/grh-lean/blob/72fead102f663d5c7a407964a13dc0cf9d977207/grh-compact/src/New/Data/Weights.lean\#L14-L16}{{\leanfont Data.\allowbreak{}nominalWeight\_mem\_Icc}}}{Weights.lean}{14--16}
\crow{nwnonneg}{\href{https://github.com/lsernau/grh-lean/blob/72fead102f663d5c7a407964a13dc0cf9d977207/grh-compact/src/New/Data/Weights.lean\#L21-L21}{{\leanfont Data.\allowbreak{}nominalWeight\_nonneg}}}{Weights.lean}{21--21}
\crow{nwadd}{\href{https://github.com/lsernau/grh-lean/blob/72fead102f663d5c7a407964a13dc0cf9d977207/grh-compact/src/New/Data/Weights.lean\#L25-L27}{{\leanfont Data.\allowbreak{}nominalWeight\_add\_le}}}{Weights.lean}{25--27}
\crow{nwlew}{\href{https://github.com/lsernau/grh-lean/blob/72fead102f663d5c7a407964a13dc0cf9d977207/grh-compact/src/New/Data/Weights.lean\#L29-L30}{{\leanfont Data.\allowbreak{}nominalWeight\_le\_w\_sub}}}{Weights.lean}{29--30}
\crow{wlenw}{\href{https://github.com/lsernau/grh-lean/blob/72fead102f663d5c7a407964a13dc0cf9d977207/grh-compact/src/New/Data/Weights.lean\#L36-L37}{{\leanfont Data.\allowbreak{}w\_le\_nominalWeight\_add}}}{Weights.lean}{36--37}
\crow{wletwob}{\href{https://github.com/lsernau/grh-lean/blob/72fead102f663d5c7a407964a13dc0cf9d977207/grh-compact/src/New/Data/Weights.lean\#L40-L41}{{\leanfont Data.\allowbreak{}w\_le\_two\_mul\_b}}}{Weights.lean}{40--41}
\crow{logw}{\href{https://github.com/lsernau/grh-lean/blob/72fead102f663d5c7a407964a13dc0cf9d977207/grh-compact/src/New/Data/Weights.lean\#L51-L54}{{\leanfont Data.\allowbreak{}log\_w\_le}}}{Weights.lean}{51--54}
\crow{chequiv}{\href{https://github.com/lsernau/grh-lean/blob/72fead102f663d5c7a407964a13dc0cf9d977207/grh-compact/src/New/Data/ChosenPlaceLevels.lean\#L18-L19}{{\leanfont Data.\allowbreak{}chosenEquiv}}}{ChosenPlaceLevels.lean}{18--19}
\crow{cardch}{\href{https://github.com/lsernau/grh-lean/blob/72fead102f663d5c7a407964a13dc0cf9d977207/grh-compact/src/New/Data/ChosenPlaceLevels.lean\#L23-L24}{{\leanfont Data.\allowbreak{}card\_chosenPlace}}}{ChosenPlaceLevels.lean}{23--24}
\crow{prodq}{\href{https://github.com/lsernau/grh-lean/blob/72fead102f663d5c7a407964a13dc0cf9d977207/grh-compact/src/New/Data/ChosenPlaceLevels.lean\#L38-L39}{{\leanfont Data.\allowbreak{}prod\_q\_chosenPlace}}}{ChosenPlaceLevels.lean}{38--39}
\crow{chstar}{\href{https://github.com/lsernau/grh-lean/blob/72fead102f663d5c7a407964a13dc0cf9d977207/grh-compact/src/New/Data/ChosenPlaceLevels.lean\#L46-L47}{{\leanfont Data.\allowbreak{}isChosen\_star\_iff}}}{ChosenPlaceLevels.lean}{46--47}
\crow{notstar}{\href{https://github.com/lsernau/grh-lean/blob/72fead102f663d5c7a407964a13dc0cf9d977207/grh-compact/src/New/Data/ChosenPlaceLevels.lean\#L81-L82}{{\leanfont Data.\allowbreak{}IsChosen.\allowbreak{}not\_star}}}{ChosenPlaceLevels.lean}{81--82}
\crow{rangech}{\href{https://github.com/lsernau/grh-lean/blob/72fead102f663d5c7a407964a13dc0cf9d977207/grh-compact/src/New/Data/ChosenPlaceLevels.lean\#L59-L60}{{\leanfont Data.\allowbreak{}range\_chosen}}}{ChosenPlaceLevels.lean}{59--60}
\crow{conjapart}{\href{https://github.com/lsernau/grh-lean/blob/72fead102f663d5c7a407964a13dc0cf9d977207/grh-compact/src/New/Data/ChosenPlaceLevels.lean\#L56-L57}{{\leanfont Data.\allowbreak{}isConjApart\_chosen}}}{ChosenPlaceLevels.lean}{56--57}
\crow{valN}{\href{https://github.com/lsernau/grh-lean/blob/72fead102f663d5c7a407964a13dc0cf9d977207/grh-compact/src/New/Data/ChosenPlaceLevels.lean\#L26-L28}{{\leanfont Data.\allowbreak{}v\_N\_eq\_exp\_neg\_one}}}{ChosenPlaceLevels.lean}{26--28}
\crow{valpi}{\href{https://github.com/lsernau/grh-lean/blob/72fead102f663d5c7a407964a13dc0cf9d977207/grh-compact/src/New/Data/ChosenPlaceLevels.lean\#L30-L32}{{\leanfont Data.\allowbreak{}intValuation\_pi}}}{ChosenPlaceLevels.lean}{30--32}
\crow{ordch}{\href{https://github.com/lsernau/grh-lean/blob/72fead102f663d5c7a407964a13dc0cf9d977207/grh-compact/src/New/Data/ChosenPlaceLevels.lean\#L36-L36}{{\leanfont Data.\allowbreak{}chosenGroupOrder}}}{ChosenPlaceLevels.lean}{36--36}
\crow{chlevel}{\href{https://github.com/lsernau/grh-lean/blob/72fead102f663d5c7a407964a13dc0cf9d977207/grh-compact/src/New/Data/ChosenPlaceLevels.lean\#L44-L44}{{\leanfont Data.\allowbreak{}chosenLevel}}}{ChosenPlaceLevels.lean}{44--44}
\crow{cardGL}{\href{https://github.com/lsernau/grh-lean/blob/72fead102f663d5c7a407964a13dc0cf9d977207/grh-compact/src/New/Data/ChosenPlaceLevels.lean\#L51-L54}{{\leanfont Data.\allowbreak{}card\_GL\_dvd\_chosenGroupOrder}}}{ChosenPlaceLevels.lean}{51--54}
\crow{princ}{\href{https://github.com/lsernau/grh-lean/blob/72fead102f663d5c7a407964a13dc0cf9d977207/grh-compact/src/New/Data/ChosenPlaceLevels.lean\#L62-L63}{{\leanfont Data.\allowbreak{}principalLevel\_le\_chosenLevel}}}{ChosenPlaceLevels.lean}{62--63}
\crow{neat}{\href{https://github.com/lsernau/grh-lean/blob/72fead102f663d5c7a407964a13dc0cf9d977207/grh-compact/src/New/Data/ChosenPlaceLevels.lean\#L67-L70}{{\leanfont Data.\allowbreak{}classAut\_eq\_bot\_of\_le\_chosenLevel}}}{ChosenPlaceLevels.lean}{67--70}
\crow{mink}{\href{https://github.com/lsernau/grh-lean/blob/72fead102f663d5c7a407964a13dc0cf9d977207/grh-compact/src/New/Data/ChosenPlaceLevels.lean\#L75-L79}{{\leanfont Data.\allowbreak{}card\_classAut\_dvd\_chosenGroupOrder}}}{ChosenPlaceLevels.lean}{75--79}
\crow{omega}{\href{https://github.com/lsernau/grh-lean/blob/72fead102f663d5c7a407964a13dc0cf9d977207/grh-compact/src/New/Data/ChosenPlaceLevels.lean\#L86-L90}{{\leanfont Data.\allowbreak{}omega\_zpow\_ne\_one}}}{ChosenPlaceLevels.lean}{86--90}
\crow{behrend}{\href{https://github.com/lsernau/grh-lean/blob/72fead102f663d5c7a407964a13dc0cf9d977207/grh-compact/src/New/Data/DataExist.lean\#L20-L36}{{\leanfont exists\_three\_term\_free}}}{DataExist.lean}{20--36}
\crow{pigeon}{\href{https://github.com/lsernau/grh-lean/blob/72fead102f663d5c7a407964a13dc0cf9d977207/grh-compact/src/New/Data/DataExist.lean\#L41-L45}{{\leanfont exists\_mul\_principal}}}{DataExist.lean}{41--45}
\crow{exdata}{\href{https://github.com/lsernau/grh-lean/blob/72fead102f663d5c7a407964a13dc0cf9d977207/grh-compact/src/New/Data/DataExist.lean\#L49-L108}{{\leanfont exists\_data}}}{DataExist.lean}{49--108}}
% the notations of the defined functions: one slot for each explicit argument of the Lean declaration
\newcommand{\twd}[1]{\operatorname{tw}(#1)}
\newcommand{\Nn}[2]{N_1(#1,#2)}
\newcommand{\KP}[1]{K(#1)}
\newcommand{\Mom}[2]{M_\omega(#1,#2)}
\newcommand{\EP}[2]{E(#1,#2)}
\newcommand{\PlP}[2]{\operatorname{Pl}(#1,#2)}
\newcommand{\pin}[3]{[#3]_{#1,#2}}
\newcommand{\pcl}[2]{\operatorname{cl}_{#1}^{#2}}
\newcommand{\OmP}[2]{\Omega(#1,#2)}
\newcommand{\CP}[1]{C(#1)}
\newcommand{\scx}[1]{x(#1)}
\newcommand{\scxo}[1]{x_1(#1)}
\newcommand{\scb}[1]{b(#1)}
\newcommand{\ED}[1]{E(#1)}
\newcommand{\PlD}[1]{\operatorname{Pl}(#1)}
\newcommand{\ND}[1]{N_1(#1)}
\newcommand{\fD}[1]{f(#1)}
\newcommand{\rdD}[1]{\operatorname{rd}(#1)}
\newcommand{\pinD}[2]{[#2]_{#1}}
\newcommand{\tops}[1]{\operatorname{tops}(#1)}
\newcommand{\gD}[1]{g(#1)}
\newcommand{\kD}[1]{k(#1)}
\newcommand{\tD}[1]{t(#1)}
\newcommand{\nD}[2]{n_{#2}(#1)}
\newcommand{\posD}[2]{\operatorname{pos}_{#2}(#1)}
\newcommand{\hhD}[1]{h_2(#1)}
\newcommand{\hD}[1]{h(#1)}
\newcommand{\rhg}[1]{h_2'(#1)}
\newcommand{\rup}[2]{\lceil #2\rceil_{#1}}
\newcommand{\aD}[1]{a(#1)}
\newcommand{\bone}[1]{b_1(#1)}
\newcommand{\btwo}[1]{b_2(#1)}
\newcommand{\wt}[3]{w_{#1,#2}(#3)}
\newcommand{\nw}[2]{w_{#1}(#2)}
\newcommand{\Ask}[3]{\operatorname{Ask}_{#1}(#2,#3)}
\newcommand{\budg}[1]{\operatorname{budget}(#1)}
\newcommand{\jet}[2]{\operatorname{jet}_{#1}(#2)}
\newcommand{\UD}[1]{\mathsf{U}(#1)}
\newcommand{\hlb}[4]{\operatorname{Bd}^{#2}_{#1}(#3;\,#4)}
\newcommand{\hlu}[1]{\operatorname{Bd}^{\top}_{#1}}
\newcommand{\hll}[1]{\operatorname{Bd}^{\bot}_{#1}}
\newcommand{\UnitMembers}{239}
\newcommand{\UnitInNote}{98}
\newcommand{\UnitProofs}{39}
\newcommand{\UnitSteps}{145}
\newcommand{\UnitAllSteps}{172}


\linespread{1.08}                       % 11 pt on 14.7 pt
\setlength{\parindent}{0pt}
\setlength{\parskip}{5.5pt plus 1.5pt minus 0.5pt}
\setlength{\abovedisplayskip}{7.5pt plus 2pt minus 3pt}
\setlength{\belowdisplayskip}{7.5pt plus 2pt minus 3pt}
\setlength{\abovedisplayshortskip}{3pt plus 2pt}
\setlength{\belowdisplayshortskip}{6pt plus 2pt minus 2pt}
\emergencystretch=1.5em
\clubpenalty=10000 \widowpenalty=10000 \displaywidowpenalty=10000
\relpenalty=7000 \binoppenalty=9000
\raggedbottom
\pagestyle{fancy}\fancyhf{}\renewcommand{\headrulewidth}{0pt}
\fancyfoot[C]{\sffamily\fontsize{7.6}{9}\selectfont\color{inkfaint}\thepage\,/\,\pageref*{LastPage}}
\setlist{topsep=2pt,partopsep=0pt,parsep=0pt,itemsep=4pt,leftmargin=2.1em,labelsep=0.55em}

\newcommand{\Q}{\mathbb{Q}}
\newcommand{\Z}{\mathbb{Z}}
\newcommand{\F}{\mathbb{F}}
\newcommand{\len}{\operatorname{len}}
\newcommand{\cok}{\operatorname{cok}}
\newcommand{\End}{\operatorname{End}}
\newcommand{\Hom}{\operatorname{Hom}}
\newcommand{\Gal}{\operatorname{Gal}}
\newcommand{\diag}{\operatorname{diag}}
\newcommand{\kw}[1]{\emph{#1}}
\newcommand{\m}[1]{\mbox{$#1$}}        % a short formula that must not be broken across lines
\newcommand{\head}[1]{\par\addvspace{17pt plus 4pt minus 2pt}\Needspace{6\baselineskip}%
  \noindent{\fontsize{13.2}{16}\selectfont\bfseries\color{ink}#1}\par\nopagebreak\vspace{-1pt}}
\newcommand{\sect}[2]{\refstepcounter{section}\label{#2}\par\addvspace{19pt plus 4pt minus 2pt}\Needspace{7\baselineskip}%
  \noindent\llap{{\fontsize{13.2}{16}\selectfont\bfseries\color{inkfaint}\thesection}\hspace{9pt}}%
  {\fontsize{13.2}{16}\selectfont\bfseries\color{ink}#1}\par\nopagebreak\vspace{-1pt}}
\newcommand{\subhead}[1]{\par\addvspace{9pt plus 2pt}\Needspace{4\baselineskip}\noindent\textbf{#1}\enspace\ignorespaces}
% where a statement lives in the Lean archive
\newcommand{\leantag}{{\sffamily\bfseries\fontsize{6.3}{8}\selectfont\color{leancol}\addfontfeatures{LetterSpace=7}LEAN}}
\newcommand{\sep}{\unskip\nobreak\hspace{5pt}{\color{inkfaint}\textperiodcentered}\hspace{5pt plus 2pt}\ignorespaces}
\newcommand{\leanline}[1]{\par\nopagebreak\vspace{-2.5pt}{\fontsize{8.9}{12.4}\selectfont\color{inksoft}\raggedright
  \noindent\llap{\leantag\hspace{9pt}}#1\par}\vspace{1pt}}
% a displayed statement in a tinted box: #1 colour, #2 text
\newcommand{\statement}[2]{\par\addvspace{7pt}\begingroup\setlength{\fboxsep}{8.5pt}\setlength{\fboxrule}{0.6pt}%
  \noindent\fcolorbox{#1!30!white}{#1!3!white}{\parbox{\dimexpr\linewidth-2\fboxsep-2\fboxrule\relax}{#2}}%
  \endgroup\par\addvspace{7pt}}
\newcommand{\nt}[2]{\par\noindent\hangindent=1.1em\hangafter=1 #1\enspace{\color{inksoft}#2}\par\vspace{1pt}}
\newcommand{\legdot}[1]{\tikz[baseline=-0.62ex]\fill[#1] (0,0) circle (2.5pt);\hspace{3pt}}
\newcommand{\crow}[4]{\csname CS@#1\endcsname & #2 & \textcolor{inksoft}{#3} & \textcolor{inksoft}{#4}\\}
\newcommand{\leanquoteline}[2]{\par\noindent\hangindent=\dimexpr#1\fontdimen2\font+2.2em\relax\hangafter=1
  \hspace*{#1\fontdimen2\font}#2}
\newcommand{\quoted}[1]{\par\vspace{3pt}{\ttfamily\fontsize{7.5}{10.4}\selectfont\color{leancol}\raggedright
  \setlength{\parskip}{0pt}\leanquote{#1}\par}\vspace{4pt}}
\newcommand{\bx}[3]{{\sffamily\bfseries\fontsize{7.3}{9.2}\selectfont\color{#1!82!black}#2}\\[1.5pt]#3}
\newcommand{\cs}[2]{\expandafter\newcommand\csname CS@#1\endcsname{#2}}

% theorems
\newtheoremstyle{plainx}{7pt plus 2pt}{7pt plus 2pt}{\itshape}{}{\bfseries}{.}{0.5em}{}
\newtheoremstyle{uprightx}{7pt plus 2pt}{7pt plus 2pt}{}{}{\bfseries}{.}{0.5em}{}
\theoremstyle{plainx}
\newtheorem{theorem}{Theorem}[section]
\newtheorem{lemma}[theorem]{Lemma}
\newtheorem{proposition}[theorem]{Proposition}
\newtheorem{corollary}[theorem]{Corollary}
\theoremstyle{uprightx}
\newtheorem{definition}[theorem]{Definition}
\newtheorem{setting}[theorem]{Setting}          % hypotheses shared by several statements; each of them cites it by number
\newtheorem{remark}[theorem]{Remark}
\newtheorem{example}[theorem]{Example}
\renewcommand{\qedsymbol}{{\setlength{\fboxsep}{0pt}\setlength{\fboxrule}{0.5pt}\fbox{\rule{0pt}{0.46em}\hspace{0.46em}}}}
% what holds throughout the note: the ambient objects only (implicit and instance arguments), never a hypothesis
\newenvironment{conventions}{\par\addvspace{5pt}\noindent\textbf{Conventions.}\enspace\ignorespaces}{\par\addvspace{3pt}}
\makeatletter
\renewenvironment{proof}[1][Proof]{\par\pushQED{\qed}\normalfont\topsep0pt\partopsep0pt
  \trivlist\item[\hskip\labelsep\itshape #1.]\ignorespaces}{\popQED\endtrivlist\@endpefalse}
\makeatother

% where each declaration of the concordance is treated: \cs{key}{place}, the place naming a \label or a lettered item
\cs{twd}{Definition~\ref{d:recall}}
\cs{Nn}{Definition~\ref{d:recall}}
\cs{KP}{Definition~\ref{d:recall}}
\cs{Mom}{Definition~\ref{d:recall}}
\cs{EP}{Definition~\ref{d:recall}}
\cs{PlP}{Definition~\ref{d:recall}}
\cs{pin}{Definition~\ref{d:recall}}
\cs{pcl}{Definition~\ref{d:recall}}
\cs{OmP}{Definition~\ref{d:recall}}
\cs{CP}{Definition~\ref{d:C}}
\cs{Data}{Definition~\ref{d:data}}
\cs{scx}{Definition~\ref{d:scales}}
\cs{scxo}{Definition~\ref{d:scales}}
\cs{scb}{Definition~\ref{d:scales}}
\cs{ED}{Definition~\ref{d:field}}
\cs{PlD}{Definition~\ref{d:field}}
\cs{ND}{Definition~\ref{d:field}}
\cs{fD}{Definition~\ref{d:field}}
\cs{cond}{Definition~\ref{d:field}}
\cs{rdD}{Definition~\ref{d:field}}
\cs{pinD}{Definition~\ref{d:field}}
\cs{tops}{Definition~\ref{d:field}}
\cs{gD}{Definition~\ref{d:field}}
\cs{kD}{Definition~\ref{d:ranks}}
\cs{tD}{Definition~\ref{d:ranks}}
\cs{nD}{Definition~\ref{d:ranks}}
\cs{posD}{Definition~\ref{d:ranks}}
\cs{hhD}{Definition~\ref{d:ranks}}
\cs{hD}{Definition~\ref{d:ranks}}
\cs{rhg}{Definition~\ref{d:ranks}}
\cs{rup}{Definition~\ref{d:order}}
\cs{aD}{Definition~\ref{d:order}}
\cs{bone}{Definition~\ref{d:order}}
\cs{btwo}{Definition~\ref{d:order}}
\cs{wt}{Definition~\ref{d:order}}
\cs{nw}{Definition~\ref{d:order}}
\cs{Ask}{Definition~\ref{d:sums}}
\cs{budg}{Definition~\ref{d:sums}}
\cs{jet}{Definition~\ref{d:sums}}
\cs{UD}{Definition~\ref{d:sums}}
\cs{hlb}{Definition~\ref{d:bd}}
\cs{hlu}{Definition~\ref{d:bd}}
\cs{hll}{Definition~\ref{d:bd}}
\cs{pos}{Lemma~\ref{l:rules}}
\cs{upabs}{Lemma~\ref{l:rules}}
\cs{upnonneg}{Lemma~\ref{l:rules}}
\cs{upnat}{Lemma~\ref{l:rules}}
\cs{lowle}{Lemma~\ref{l:rules}}
\cs{icc}{Lemma~\ref{l:icc}}
\cs{ltBd}{Theorem B}
\cs{leBd}{Theorem B}
\cs{natlt}{Theorem B}
\cs{natle}{Theorem B}
\cs{natcast}{Proposition~\ref{p:table}}
\cs{bdmul}{Proposition~\ref{p:table}}
\cs{bdadd}{Proposition~\ref{p:table}}
\cs{bdf}{Proposition~\ref{p:table}}
\cs{cge}{Lemma~\ref{l:cge}}
\cs{round}{Lemma~\ref{l:round}}
\cs{rodd}{Lemma~\ref{l:ranks}}
\cs{meven}{Lemma~\ref{l:ranks}}
\cs{lem}{Lemma~\ref{l:ranks}}
\cs{radd}{Lemma~\ref{l:ranks}}
\cs{kler}{Lemma~\ref{l:ranks}}
\cs{hundred}{Lemma~\ref{l:ranks}}
\cs{kpos}{Lemma~\ref{l:ranks}}
\cs{nmono}{Lemma~\ref{l:ranks}}
\cs{keven}{Lemma~\ref{l:ranks}}
\cs{todd}{Lemma~\ref{l:ranks}}
\cs{rsix}{Lemma~\ref{l:ranks}}
\cs{kth}{Lemma~\ref{l:ranks}}
\cs{clthick}{Lemma~\ref{l:thick}}
\cs{nwicc}{Lemma~\ref{l:weights}}
\cs{nwnonneg}{Lemma~\ref{l:weights}}
\cs{nwadd}{Lemma~\ref{l:weights}}
\cs{nwlew}{Lemma~\ref{l:weights}}
\cs{wlenw}{Lemma~\ref{l:weights}}
\cs{wletwob}{Lemma~\ref{l:weights}}
\cs{logw}{Lemma~\ref{l:weights}}
\cs{chequiv}{Lemma~\ref{l:chosen}}
\cs{cardch}{Lemma~\ref{l:chosen}}
\cs{prodq}{Lemma~\ref{l:chosen}}
\cs{chstar}{Lemma~\ref{l:chosen}}
\cs{notstar}{Lemma~\ref{l:chosen}}
\cs{rangech}{Lemma~\ref{l:chosen}}
\cs{conjapart}{Lemma~\ref{l:chosen}}
\cs{valN}{Lemma~\ref{l:level}}
\cs{valpi}{Lemma~\ref{l:level}}
\cs{ordch}{Definition~\ref{d:chosen}}
\cs{chlevel}{Definition~\ref{d:chosen}}
\cs{cardGL}{Lemma~\ref{l:GL}}
\cs{princ}{Lemma~\ref{l:GL}}
\cs{neat}{Proposition~\ref{p:neat}}
\cs{mink}{Proposition~\ref{p:neat}}
\cs{omega}{Lemma~\ref{l:omega}}
\cs{behrend}{Proposition~\ref{p:behrend}}
\cs{pigeon}{Lemma~\ref{l:pigeon}}
\cs{exdata}{Theorem A}

% ------------------------------------------------------------------ macros of this note that are not notations of
% defined functions (those are made by the build, in data.tex, from the spec)
\newcommand{\OO}{\mathcal{O}}
\newcommand{\CC}{\mathbb{C}}
\newcommand{\RR}{\mathbb{R}}
\newcommand{\NN}{\mathbb{N}}
\newcommand{\Dif}{\mathfrak{D}}
\newcommand{\Cy}[1]{\Psi_{#1}}                 % the cyclotomic polynomial
\newcommand{\chiK}[1]{\chi_{\KP{#1}}}         % the quadratic character of K(P)
\newcommand{\emb}[1]{\iota_{#1}}              % the chosen embedding at an infinite place
\newcommand{\cj}[1]{\overline{#1}}
\newcommand{\ord}{\operatorname{ord}}
\newcommand{\Nm}{\operatorname{N}}
\newcommand{\Aut}{\operatorname{Aut}}
\newcommand{\GL}{\operatorname{GL}}
% components of the data D (fields of the record; not declarations of their own)
\newcommand{\Dm}[1]{m(#1)}
\newcommand{\Dr}[1]{r(#1)}
\newcommand{\Dmu}[1]{\mu(#1)}
\newcommand{\thick}[1]{\operatorname{thick}(#1)}
\newcommand{\Dg}[2]{g_{#1}(#2)}
\newcommand{\Dpig}[1]{\pi_g(#1)}
\newcommand{\ordch}[1]{\operatorname{ord}_{\mathrm{ch}}(#1)}
\newcommand{\Kch}[2]{K_{\mathrm{ch}}(#1,#2)}
\newcommand{\rdE}[2]{\operatorname{rd}_{\EP{#1}{#2}}}

\newcommand{\lvl}[1]{N(#1)}
\newcommand{\piN}[1]{\pi_N(#1)}
\newcommand{\Gf}[2]{G_{\mathrm{f}}(#1,#2)}
\newcommand{\Kz}[2]{K_0(#1,#2)}
\newcommand{\Gam}[2]{\Gamma(#1,#2)}
\newcommand{\Ch}[2]{\operatorname{Ch}_{#1}(#2)}
\newcommand{\clno}[1]{\operatorname{h}(\KP{#1})}
\newcommand{\lx}[1]{\log\scx{#1}}
\newcommand{\lb}[1]{\log\scb{#1}}
\newcommand{\lxo}[1]{\log\scxo{#1}}
\newcommand{\llb}[1]{\log\log\scb{#1}}
\begin{document}
\color{ink}

{\raggedright
{\sffamily\bfseries\fontsize{7.2}{9}\selectfont\color{leancol}\addfontfeatures{LetterSpace=8}READBACK NOTE 6 \textperiodcentered{} CHAPTER III, SECOND PART\par}
\vspace{5pt}
{\fontsize{21.5}{25.5}\selectfont\bfseries Admissible data: scales, ranks and sizes, and their existence\par}
\vspace{5pt}
{\fontsize{12.4}{16}\selectfont\itshape\color{inksoft} One unit of the Lean artifact \ArtifactName, read back into
mathematics\par}}
\vspace{8pt}
{\color{rulegrey}\hrule height 0.6pt}
\vspace{7pt}
\thispagestyle{empty}

\textbf{Summary.} The whole argument is run on one record of numbers and places, the \emph{admissible data} $D$:
three scales $\scx{D}$, $\scxo{D}$, $\scb{D}$ with $\lb{D}\ll\scx{D}\ll\scxo{D}\ll\scb{D}$, two ranks
$\Dm{D}\approx\sqrt{\scx{D}}$ and $\Dr{D}\approx\scx{D}^5$, a set of \emph{thick} infinite places of the CM field
$\ED{D}$ that satisfies no multiplicative relation with small exponents, and two auxiliary primes. This note states
the record, names what the later chapters compute from it (ranks, the order $\aD{D}$, weights, the prime sum, the
budget), and describes the device by which the formal text compares sizes: every quantity carries an upper and a
lower exponent, and two quantities are compared by their exponents (Theorem~B). It then proves that admissible data
exist at every large scale (Theorem~A), by a weak form of the Siegel--Walfisz theorem, a pigeonhole in the class
group and a construction of Behrend type.

\statement{leancol}{\fontsize{9.6}{13.2}\selectfont\textbf{What this note is.} The artifact \ArtifactName{} is a
body of Lean code whose final theorems state the Riemann hypothesis and its analogue for Dirichlet $L$-functions.
That is a claim: its authors report that Lean's kernel accepts the code, a replay of that check by the owner of
this write-up is recorded elsewhere, and nothing was compiled for this note; no mathematician has read the proof
through. This note treats the second part of one unit (the folder of the data of the proof): \UnitInNote{} of its
\UnitMembers{} declarations, among them nine definitions recalled from the first part, the note \emph{The parameters
and the CM field}. Every statement was compared, argument by argument, with the Lean text. Proofs are complete,
with two exceptions that are said where they occur: results of other units and of the background library are
quoted as \emph{Inputs}, and the comparisons of exponents that the formal text leaves to a decision procedure are
carried out here by hand only in the cases listed in Section~\ref{s:notes}. The Lean text read is \PinText;
Appendix~A is the concordance of \ConcDecls{} declarations.}

\sect{The statements}{s:intro}
\statement{arith}{\phantomsection\label{t:A}\textbf{Theorem A (admissible data exist).}\enspace Let $P=(N_0,d)$ be
parameters and $\varepsilon$ a real number with $0<\varepsilon/6$ and $\varepsilon<1/100$. Then for every
sufficiently large real number $y$ there are admissible data $D$ for $P$ (Definition~\ref{d:data}) such that
\[\scx{D}=y,\qquad \lxo{D}=y^{\varepsilon/6},\qquad (10^6\log y)^2\le\lxo{D},\]
\[\llb{D}=\varepsilon\log y,\qquad \twd{\Dmu{D}}/\lb{D}\le\#\thick{D}.\]}

By Theorem~A of the first part, $\twd{\Dmu{D}}$ is the number of the places of class $1$ modulo $N_0$, the set
$\tops{D}$. So the last inequality says that the thick places are at least as many as the tops divided by
$\lb{D}$.

\statement{arith}{\phantomsection\label{t:B}\textbf{Theorem B (comparison by exponents).}\enspace Let $D$ be admissible
data for parameters $P=(N_0,d)$, and let $l$, $l'$ be real numbers with $l<l'$.
\begin{enumerate}[label=(\alph*)]
\item If $F$, $G$ are real numbers with $\hlu{D}(F;\,l)$ and $\hll{D}(G;\,l')$, that is $|F|\le e^{l}$ and
$e^{l'}\le G$, then $F<G$, and so $F\le G$.
\item If $A$, $B$ are natural numbers with $\hlu{D}(A;\,l)$ and $\hll{D}(B;\,l')$, then $A<B$, and so $A\le B$.
\end{enumerate}}

The statement is immediate. What it stands for is a method. In the formal text the two bounds are found
mechanically from the shape of the expressions $F$ and $G$ (Proposition~\ref{p:table} and Table~1), and the
inequality $l<l'$ between two linear expressions in $\lx{D}$, $\lxo{D}$, $\lb{D}$, $\llb{D}$ is found from the
conditions on the scales. All comparisons of sizes of the later chapters are made in this way.

\sect{The setting}{s:setting}
\begin{conventions}\phantomsection\label{c:conventions}
$\log$ is the natural logarithm and $e$ its base; for natural numbers, subtraction stops at $0$ and $\lfloor\ \rfloor$ is
the floor, as in the formal text. A \emph{finite place} of a number field is a nonzero prime ideal of its ring of
integers, written like the ideal, with norm $\Nm(w)$; an \emph{infinite place} is an archimedean absolute value.
The bar is complex conjugation on a CM field and on its finite places. An index set $I$ is a finite set.
\end{conventions}

\begin{definition}[Recalled from the first part]\label{d:recall}
\emph{Parameters} are a pair $P=(N_0,d)$ of natural numbers with $18480\mid N_0$, $d$ prime, $N_0\mid2(d+1)$ and
$N_0\nmid d+1$. For a natural number $\mu$: $\twd{\mu}=3^{\lfloor\log_3\mu\rfloor+5}2^{\mu}$,
$\Nn{P}{\mu}=N_0\,\twd{\mu}$ and $\Mom{P}{\mu}=(N_0d)^{\mu+6}$; $\KP{P}=\Q(\sqrt{-d})\subset\CC$, with quadratic
character $\chiK{P}$; $\EP{P}{\mu}$ is $\KP{P}$ with the roots of unity of order $\Nn{P}{\mu}$ adjoined, a CM field;
$\PlP{P}{\mu}$ is the set of its infinite places, a commutative group isomorphic to $(\Z/\Nn{P}{\mu})^\times$; for
an integer $\nu$, $\pin{P}{\mu}{\nu}\in\PlP{P}{\mu}$ is the place whose class is $\nu$ if $\nu$ is prime to
$\Nn{P}{\mu}$, and the unit element otherwise; for a natural number $m$ with $m\mid\Nn{P}{\mu}$,
$\pcl{P}{m\mid\Nn{P}{\mu}}\colon\PlP{P}{\mu}\to(\Z/m)^\times$ is the class modulo $m$; and
$\OmP{P}{\mu}=\exp(\mathrm{cs}(d))\cdot2\pi/\mathrm{rd}$, with $\mathrm{rd}$ the root discriminant of $\EP{P}{\mu}$
and $\mathrm{cs}(d)$ Euler's constant minus the real part of the logarithmic derivative at $1$ of the $L$-function
of $\chiK{P}$.
\end{definition}
\leanline{\lean{twd}\sep\lean{Nn}\sep\lean{Mom}\sep\lean{KP}\sep\lean{EP}\sep\lean{PlP}\sep\lean{pin}\sep\lean{pcl}\sep
\lean{OmP}}

\begin{definition}[The constant of the parameters]\label{d:C}
For parameters $P=(N_0,d)$,
\[\CP{P}=C_{\mathrm{jet}}+(2d+6)+\Bigl(N_0+|\mathrm{cs}(d)|+\sum_{x}\bigl|\log|\Cy{N_0}'(x)|\bigr|\Bigr),\]
where $x$ runs over the complex primitive roots of unity of order $N_0$, $\Cy{N_0}$ is the cyclotomic polynomial of
that order, and $C_{\mathrm{jet}}$ is a natural number chosen from an existence statement of unit~150 (there it bounds the size
of the conjugates of certain algebraic numbers of the Weil representation; this unit uses nothing about it). What
this unit uses of $\CP{P}$ is Lemma~\ref{l:cge}.
\end{definition}
\leanline{\lean{CP}}

\begin{definition}[Admissible data]\label{d:data}
Let $P=(N_0,d)$ be parameters. \emph{Admissible data for $P$} are a record $D$ with the following components and
properties. Components: three real numbers, the logarithms $\lx{D}$, $\lxo{D}$, $\lb{D}$ of the scales of
Definition~\ref{d:scales}; natural numbers $\Dmu{D}$, $\Dm{D}$, $\Dr{D}$; a finite set $\thick{D}$ of infinite
places of $\EP{P}{\Dmu{D}}$, the \emph{thick places}; natural numbers $\Dg{0}{D}$, $\Dg{1}{D}$; and an integer
$\Dpig{D}$ of $\KP{P}$. Properties:
\begin{enumerate}[label=(D\arabic*)]
\item $10^{12}\,\CP{P}\le\llb{D}$, \ $100\,\llb{D}\le\lx{D}$, \ $10^6\,\lx{D}\le\lxo{D}$, \ $10^6\,\lxo{D}\le\lb{D}$ and
$6\log\lxo{D}\le\llb{D}$.
\item $\Dmu{D}\le32\,\lx{D}$.
\item $\sqrt{\scx{D}}\le\Dm{D}\le2\sqrt{\scx{D}}$, and there is a natural number $n$ with $\Dm{D}=2^{n+1}$ and
$\Dmu{D}=22(n+1)$.
\item $\scx{D}^5\le\Dr{D}\le3\,\scx{D}^5$, and there is a natural number $n$ with $\Dr{D}=3^{n+1}$.
\item For every natural number $m$ with $m\mid\Nn{P}{\Dmu{D}}$ that divides
$N_0\cdot3^{\lfloor\log_3\Dmu{D}\rfloor+5}$, every thick place has class $-1$ modulo $m$:
$\pcl{P}{m\mid\Nn{P}{\Dmu{D}}}(\sigma)=-1$ for all $\sigma\in\thick{D}$.
\item Let $\psi=\pcl{P}{2^{\Dmu{D}+1}\mid\Nn{P}{\Dmu{D}}}$ be the class modulo $2^{\Dmu{D}+1}$. For all natural numbers $i$, $j$
with $0<i<j\le324\,\Dmu{D}$ and all $\sigma,\tau,\rho\in\thick{D}$: if $\psi(\tau/\sigma)^j=\psi(\rho/\sigma)^i$, then
$\psi(\tau/\sigma)=1$.
\item $\Dg{0}{D}\ne\Dg{1}{D}$; and for $i\in\{0,1\}$: $\Dg{i}{D}$ is a prime number,
$\Dg{i}{D}\equiv1\bmod\Nn{P}{\Dmu{D}}\,d$, the square of $\Mom{P}{\Dmu{D}}$ divides $\Dg{i}{D}-1$, and
$\scb{D}^{10}\le\Dg{i}{D}\le\scb{D}^{11}$.
\item $\Dpig{D}\cdot\cj{\Dpig{D}}=\Dg{0}{D}\,\Dg{1}{D}$.
\item For every real number $y$ with $100\,\lb{D}\le y\le2^{\Dmu{D}}$ there is a prime number $p$ with
$p\equiv1$ modulo the square of $\Mom{P}{\Dmu{D}}$ and $y\le\log p\le y+\sqrt{\lb{D}}$.
\end{enumerate}
\end{definition}
\leanline{\lean{Data}}

\begin{definition}[The scales]\label{d:scales}
For admissible data $D$ for parameters $P=(N_0,d)$, the three \emph{scales} $\scx{D}$, $\scxo{D}$, $\scb{D}$ are the
exponentials of the three real numbers that $D$ stores. In terms of them (D1) reads
\[\exp\bigl(10^{12}\CP{P}\bigr)\le\lb{D},\qquad(\lb{D})^{100}\le\scx{D},\qquad\scx{D}^{10^6}\le\scxo{D},\]
\[\scxo{D}^{10^6}\le\scb{D},\qquad(\lxo{D})^{6}\le\lb{D}.\]
\end{definition}
\leanline{\lean{scx}\sep\lean{scxo}\sep\lean{scb}}

\begin{definition}[The field and its places]\label{d:field}
Let $D$ be admissible data for parameters $P=(N_0,d)$. Put $\ND{D}=\Nn{P}{\Dmu{D}}$, $\ED{D}=\EP{P}{\Dmu{D}}$ and
$\PlD{D}=\PlP{P}{\Dmu{D}}$; let $\fD{D}$ be the number of elements of $\PlD{D}$ and $\rdD{D}$ the root discriminant
of $\ED{D}$. The formal text also names the product $\ND{D}\cdot d$, the order of a cyclotomic field that contains
$\ED{D}$. For an integer $j$, $\pinD{D}{j}=\pin{P}{\Dmu{D}}{j}$. The \emph{tops} are the places of class $1$ modulo
$N_0$,
\[\tops{D}=\bigl\{\tau\in\PlD{D}:\ \pcl{P}{N_0\mid\Nn{P}{\Dmu{D}}}(\tau)=1\bigr\},\]
and $\gD{D}=\Dg{0}{D}\,\Dg{1}{D}$.
\end{definition}
\leanline{\lean{ND}\sep\lean{ED}\sep\lean{PlD}\sep\lean{fD}\sep\lean{rdD}\sep\lean{cond}\sep\lean{pinD}\sep\lean{tops}\sep
\lean{gD}}

\begin{definition}[Ranks and gaps]\label{d:ranks}
Let $D$ be admissible data for parameters $P=(N_0,d)$, and $j$ a natural number. Put
\[\kD{D}=(\Dr{D}+1)-\Dm{D},\qquad\tD{D}=5\,\Dr{D},\]
\[\nD{D}{j}=\kD{D}+j\,(\tD{D}+1),\qquad\posD{D}{j}=\kD{D}+j-1,\]
\[\hhD{D}=\Dr{D}^{8},\qquad\hD{D}=2\,\hhD{D},\qquad\rhg{D}=\hhD{D}-\lfloor(\tD{D}+1)/2\rfloor .\]
\end{definition}
\leanline{\lean{kD}\sep\lean{tD}\sep\lean{nD}\sep\lean{posD}\sep\lean{hhD}\sep\lean{hD}\sep\lean{rhg}}

\begin{definition}[The order and the weights]\label{d:order}
Let $D$ be admissible data for parameters $P=(N_0,d)$, $y$ a real number, $j$ a natural number and
$\sigma\in\PlD{D}$. $\rup{D}{y}$ is the least natural multiple of $\ND{D}$ that is at least $y$ (it is $0$ for
$y\le0$). Put
\[\aD{D}=\rup{D}{\scx{D}\,\scb{D}},\qquad\bone{D}=\rup{D}{\aD{D}/\scx{D}},\qquad\btwo{D}=\rup{D}{\aD{D}/\scxo{D}} .\]
The \emph{nominal weight} $\nw{D}{\sigma}$ is $\aD{D}/\scx{D}$ if $\sigma\in\thick{D}$ and $\aD{D}/\scxo{D}$ otherwise.
The \emph{weight} $\wt{D}{j}{\sigma}$ is $\bone{D}+j\,\hD{D}+\hhD{D}$ if $\sigma\in\thick{D}$ and
$\btwo{D}+j\,\hD{D}+\hhD{D}$ otherwise.
\end{definition}
\leanline{\lean{rup}\sep\lean{aD}\sep\lean{bone}\sep\lean{btwo}\sep\lean{nw}\sep\lean{wt}}

\begin{definition}[The prime sum, the budget, the cost of one direction]\label{d:sums}
Let $D$ be admissible data for parameters $P=(N_0,d)$, $y$ a real number, $j$ a natural number and
$\sigma\in\PlD{D}$. With $\lambda_{d,1}(\nu)=\Lambda(\nu)\,(1-\chiK{P}(\nu))$, where $\Lambda$ is von Mangoldt's
function (a definition of unit~151), and $(\ )_+$ the positive part,
\[\Ask{D}{y}{\sigma}=\sum_{\nu=1}^{\aD{D}}\lambda_{d,1}(\nu)\Bigl(\frac y\nu-\nw{D}{\pinD{D}{-\nu}\cdot\sigma}\Bigr)_{+},
\qquad\budg{D}=\fD{D}\,\aD{D}\,(\lb{D})^{3},\]
\[\jet{D}{j}=\log(j!)-j\log\frac{\OmP{P}{\Dmu{D}}}{\gD{D}},\qquad
\UD{D}=\sum_{\sigma\in\PlD{D}}\Bigl(\jet{D}{\wt{D}{0}{\sigma}}-\Ask{D}{\nw{D}{\sigma}}{\sigma}\Bigr).\]
\end{definition}
\leanline{\lean{Ask}\sep\lean{budg}\sep\lean{jet}\sep\lean{UD}}

\textbf{Why these choices.} This paragraph is motivation, taken from the comments of the formal text and from the
places where later units use the data; nothing in it is a statement of this unit. The one-sided bound for primes is
proved at the scale $\scx{D}$, and the saving is $1/\Dm{D}$, about $\scx{D}^{-1/2}$: $\Dm{D}$ is the number of theta
lifts. At the much larger scale $\scxo{D}$ a second remainder enters, which Chapter~II makes harmless by moving
$\scxo{D}$ inside a window; the conditions leave room for that. The scale $\scb{D}$ is the size of the weights, which
are orders of vanishing: $\aD{D}\approx\scx{D}\scb{D}$ at the tops and about $\scb{D}$, or $\scb{D}\scx{D}/\scxo{D}$,
at the other places, according to whether the place is thick. All losses of the argument are powers of
$\lb{D}$, the unit of the budget, and $(\lb{D})^{100}\le\scx{D}$ makes them small against the saving. The ranks grow
in $2\Dm{D}$ steps of $\tD{D}+1$ from $\kD{D}$; $\Dm{D}$ is a power of $2$ and $\Dr{D}$ a power of $3$ so that the
parities of Lemma~\ref{l:ranks} hold and $2^{\Dmu{D}}=\Dm{D}^{22}$ is about $\scx{D}^{11}$. The thick places all have
class $-1$ modulo $N_0$, so for a top $\tau$ the place $\pinD{D}{-\nu}\cdot\tau$ can be thick only for
$\nu\equiv1\bmod N_0$: this is how the progression $1$ modulo $N_0$ enters the prime sum, with the number of thick
places as its coefficient, and why Theorem~A must provide many of them. Condition (D6) says that the thick places,
seen in the $2$-part of the group of places, satisfy no relation $\psi(\tau/\sigma)^j=\psi(\rho/\sigma)^i$ with
small exponents; a later unit applies it to $\tau=\pinD{D}{-p}^{\,i}\sigma$ and $\rho=\pinD{D}{-p}^{\,j}\sigma$ for a
prime $p$, so that the powers of one Frobenius class carry a thick place to at most one other thick place. The
primes $\Dg{0}{D}$, $\Dg{1}{D}$ split completely in $\ED{D}$; they carry an auxiliary level
(Section~\ref{s:chosen}), and (D9) is used in unit~141 to choose a prime level $\lvl{D}$ of prescribed size.

\sect{The calculus of sizes}{s:sizes}
\begin{definition}[Bounds on the logarithmic scale]\label{d:bd}
Let $D$ be admissible data for parameters $P=(N_0,d)$, let $F$ and $l$ be real numbers and $\ast$ one of the two
signs $\top$, $\bot$ (a truth value in the formal text). The statement $\hlb{D}{\ast}{F}{l}$ means $|F|\le e^{l}$ if
$\ast=\top$, and $e^{l}\le F$ if $\ast=\bot$. The two cases are also written $\hlu{D}(F;\,l)$ for
$\hlb{D}{\top}{F}{l}$ and $\hll{D}(F;\,l)$ for $\hlb{D}{\bot}{F}{l}$; in the formal text $\hlu{D}$ and $\hll{D}$
are the relation with its sign fixed. The data $D$ do not occur in the meaning: they are an index, which lets the
rules below name the quantities of $D$. We call $l$ an \emph{upper exponent}, respectively a \emph{lower exponent},
of $F$.
\end{definition}
\leanline{\lean{hlb}\sep\lean{hlu}\sep\lean{hll}}

\begin{lemma}[First rules]\label{l:rules}
Let $D$ be admissible data for parameters $P=(N_0,d)$, let $F$, $G$, $l$, $l'$ be real numbers and $A$, $B$ natural
numbers.
\begin{enumerate}[label=(\alph*)]
\item If $\hll{D}(F;\,l)$, then $0<F$.
\item If $\hlu{D}(F;\,l)$ and $|G|\le F$, then $\hlu{D}(G;\,l)$. If $\hlu{D}(F;\,l)$, $G\le F$ and $0\le G$, then
$\hlu{D}(G;\,l)$.
\item If $\hlu{D}(A;\,l)$ and $B\le A$, then $\hlu{D}(B;\,l)$.
\item If $\hll{D}(G;\,l')$ and $G\le F$, then $\hll{D}(F;\,l')$.
\end{enumerate}
\end{lemma}
\begin{proof}
Each is a chain of two inequalities, with $e^{l}>0$ in (a).
\end{proof}
\leanline{\lean{pos}\sep\lean{upabs}\sep\lean{upnonneg}\sep\lean{upnat}\sep\lean{lowle}}

\begin{lemma}[The interval rule]\label{l:icc}
Let $D$ be admissible data for parameters $P=(N_0,d)$, let $F$, $G$, $H$, $l$, $l'$, $l''$ be real numbers with
$\hll{D}(G;\,l)$, $\hlu{D}(H;\,l')$ and $G\le F\le H$, and let $\ast$ be one of the two signs. Assume $l'\le l''$ or
$l''\le l$, according to whether $\ast$ is $\top$ or $\bot$. Then $\hlb{D}{\ast}{F}{l''}$.
\end{lemma}
\begin{proof}
$0<e^{l}\le G\le F\le H\le e^{l'}$, so $|F|\le e^{l'}\le e^{l''}$ in the first case and $e^{l''}\le e^{l}\le F$ in the
second.
\end{proof}
\leanline{\lean{icc}. In the formal text $G$ and $H$ are $F$ by default, and the two last hypotheses are found by
automation.}

\begin{proof}[Proof of Theorem B]
$F\le|F|\le e^{l}<e^{l'}\le G$; in (b) the natural numbers are compared as real numbers.
\end{proof}
\leanline{\lean{ltBd}\sep\lean{leBd}\sep\lean{natlt}\sep\lean{natle}}

\begin{proposition}[The rules, and the exponents of the quantities of the data]\label{p:table}
Let $D$ be admissible data for parameters $P=(N_0,d)$. Let $F$, $G$ be real numbers with upper exponents $l$, $l'$,
or with lower exponents $l$, $l'$; in the second case they are positive. Then, in the same sense:
\begin{enumerate}[label=(\alph*)]
\item $F\,G$ has the exponent $l+l'$, and $F^{A}$ the exponent $A\,l$ for every natural number $A$;
\item $F+G$ has the upper exponent $\max(l,l')+1$, respectively the lower exponent $\max(l,l')$;
\item $\sqrt F$ has the exponent $l/2$; and if $F$ has the lower exponent $l$ and the upper exponent $l'$, then $1/F$
has the upper exponent $-l$ and the lower exponent $-l'$;
\item a natural number $A\ne0$ has the upper exponent $A$ and the lower exponent $0$, and $e^{l}$ has both
exponents $l$.
\end{enumerate}
Moreover every line of Table~1 holds: the quantity in the first column has the upper exponent in the second column
and, where one is given, the lower exponent in the third.
\end{proposition}

\begin{center}\fontsize{9}{12.4}\selectfont
\begin{tabular}{@{}lll@{}}
\toprule
quantity & upper exponent & lower exponent\\\midrule
$\lb{D}$ & $\llb{D}$ & $\llb{D}$\\
$\lxo{D}$ & $\log\lxo{D}$ & $\log\lxo{D}$\\
$\Dmu{D}$, \ $\CP{P}$ & $\log\lxo{D}$, \ $\llb{D}$ & \\
$\Dm{D}$ & $\tfrac12\lx{D}+2$ & $\tfrac12\lx{D}$\\
$\Dr{D}$; \ $\tD{D}$ & $5\lx{D}+3$; \ $5\lx{D}+8$ & $5\lx{D}$\\
$\kD{D}$ & $5\lx{D}+3$ & $5\lx{D}-2$\\
$\nD{D}{\Dm{D}}$ & $5.5\,\lx{D}+12$ & $5.5\,\lx{D}$\\
$\hhD{D}$; \ $\hD{D}$ & $40\lx{D}+25$; \ $40\lx{D}+27$ & $40\lx{D}$\\
$\twd{\Dmu{D}}$ & $11\lx{D}+\tfrac15\llb{D}$ & $11\lx{D}$\\
$\fD{D}$ & $11\lx{D}+\tfrac15\llb{D}+\CP{P}$ & $11\lx{D}$\\
$\Dg{i}{D}$; \ $\gD{D}$ & $11\lb{D}$; \ $22\lb{D}$ & $10\lb{D}$; \ $20\lb{D}$\\
$\aD{D}$ & $\lb{D}+\lx{D}+1$ & $\lb{D}+\lx{D}$\\
$\bone{D}$ & $\lb{D}+2$ & $\lb{D}-2$\\
$\btwo{D}$ & $\lb{D}-\lxo{D}+\lx{D}+2$ & $\lb{D}-\lxo{D}+\lx{D}$\\
$\nw{D}{\sigma}$ & $\lb{D}+1$ & $\lb{D}-\lxo{D}$\\
$\wt{D}{0}{\sigma}$ & $\lb{D}+2$ & that of $\nw{D}{\sigma}$, for every $\wt{D}{j}{\sigma}$\\
$\Mom{P}{\Dmu{D}}$ & $\tfrac19\lb{D}$ & \\
$\log\ordch{D}$ & $32\lx{D}$ & \\
\bottomrule
\end{tabular}

\smallskip{\fontsize{8.9}{12}\selectfont\color{inksoft}Table 1. Exponents recorded in the formal text; $\sigma\in\PlD{D}$,
$i\in\{0,1\}$, $j$ a natural number; $\ordch{D}$ is defined in Definition~\ref{d:chosen}.}
\end{center}

\begin{proof}
(a) to (d) are the rules $e^{l}e^{l'}=e^{l+l'}$, $2<e$, and $A\le e^{A}$. A line of the table follows from the
interval rule, Lemma~\ref{l:icc}, applied to the bounds that the definitions give, and from (a) to (d). For example
(D3) gives $e^{\lx{D}/2}\le\Dm{D}\le2\,e^{\lx{D}/2}$, and $2\le e^{2}$. For $\twd{\Dmu{D}}$: the factor
$2^{\Dmu{D}}=\Dm{D}^{22}$ has the exponents $11\lx{D}+44$ and $11\lx{D}$; the factor $3^{\lfloor\log_3\Dmu{D}\rfloor}$
is at most $\Dmu{D}+1$ and $3^5$ has the upper exponent $15$, so an upper exponent is
$11\lx{D}+\log\lxo{D}+60$, which is below the one in the table because $\log\lxo{D}\le\frac16\llb{D}$ and
$\llb{D}\ge10^{12}\,\CP{P}$ with $\CP{P}\ge64$ (Lemma~\ref{l:cge}). The line for $\fD{D}$ follows from this one by
$\fD{D}=\twd{\Dmu{D}}\,\varphi(N_0)$ (Theorem~A of the first part) and $\varphi(N_0)\le\CP{P}$. For $\aD{D}$ use
Lemma~\ref{l:round}: $\scx{D}\scb{D}\le\aD{D}\le\scx{D}\scb{D}+\ND{D}$. The remaining lines are verified in the same
way; we did not redo all of them by hand (Section~\ref{s:notes}).
\end{proof}
\leanline{\lean{bdmul}\sep\lean{bdadd}\sep\lean{natcast}\sep\lean{bdf}; the other rules and lines are instances without
names.}

The formal text has a few more rules of the same kind, used rarely: for the natural ceiling, the absolute value, a
difference, a finite sum, and for upper bounds of logarithms ($\log F$, $\log(FG)$, $\log F^{A}$, $\log A!$).

\begin{lemma}\label{l:cge}
Let $P=(N_0,d)$ be parameters. Then $64\le\CP{P}$, $N_0\le\CP{P}$, $d\le\CP{P}$ and $\varphi(N_0)\le\CP{P}$, with
$\varphi$ Euler's function.
\end{lemma}
\begin{proof}
All summands in Definition~\ref{d:C} are nonnegative, $64<d$, and $\varphi(N_0)\le N_0$.
\end{proof}
\leanline{\lean{cge}}

\begin{lemma}[Rounding up]\label{l:round}
Let $D$ be admissible data for parameters $P=(N_0,d)$, and let $F$, $l$ be real numbers with $\hll{D}(F;\,l)$. Then
$F\le\rup{D}{F}\le F+\ND{D}$.
\end{lemma}
\begin{proof}
$F>0$, and the least integer $n\ge F/\ND{D}$ satisfies $n<F/\ND{D}+1$; multiply by $\ND{D}$.
\end{proof}
\leanline{\lean{round}}

\begin{lemma}[Ranks]\label{l:ranks}
Let $D$ be admissible data for parameters $P=(N_0,d)$. Then:
\begin{enumerate}[label=(\alph*)]
\item $\Dr{D}$ is odd, $\Dm{D}$ is even, $\kD{D}$ is even and $\tD{D}$ is odd;
\item $10000\le\Dm{D}$; \ $\Dr{D}+1=\kD{D}+\Dm{D}$; \ $\kD{D}\le\Dr{D}$; \ $100\,\Dm{D}\le\kD{D}$; \ $0<\kD{D}$;
\item $\Dr{D}+6\le\tD{D}$; \ $\kD{D}+\tD{D}+1\le\hhD{D}$; and $j\mapsto\nD{D}{j}$ is monotone.
\end{enumerate}
\end{lemma}
\begin{proof}
By (D3) and (D4), $\Dm{D}$ is a power of $2$ and $\Dr{D}$ a power of $3$, both with positive exponent. From Table~1
and Theorem~B: $10000<\Dm{D}$, because $10000<\frac12\lx{D}$ by (D1) and Lemma~\ref{l:cge}; and
$2\,\Dm{D}\le\Dr{D}$, because $\frac12\lx{D}+4<5\lx{D}$. So the subtraction in $\kD{D}$ is exact,
$\Dr{D}/2\le\kD{D}\le\Dr{D}$, and $100\,\Dm{D}\le\kD{D}$ because $\frac12\lx{D}+102<5\lx{D}-2$. The parities follow from
$\kD{D}=\Dr{D}+1-\Dm{D}$ and $\tD{D}=5\,\Dr{D}$. Finally $\Dr{D}\ge3$ gives $\Dr{D}+6\le5\,\Dr{D}$ and
$\kD{D}+\tD{D}+1\le6\,\Dr{D}+1\le\Dr{D}^8$.
\end{proof}
\leanline{\lean{rodd}\sep\lean{meven}\sep\lean{keven}\sep\lean{todd}\sep\lean{lem}\sep\lean{radd}\sep\lean{kler}\sep
\lean{hundred}\sep\lean{kpos}\sep\lean{rsix}\sep\lean{kth}\sep\lean{nmono}}

\begin{lemma}\label{l:thick}
Let $D$ be admissible data for parameters $P=(N_0,d)$, and $\sigma\in\PlD{D}$ with $\sigma\in\thick{D}$. Then
$\pcl{P}{N_0\mid\Nn{P}{\Dmu{D}}}(\sigma)=-1$.
\end{lemma}
\begin{proof}
This is (D5) for the divisor $N_0$.
\end{proof}
\leanline{\lean{clthick}}

\sect{Weights}{s:weights}
\begin{lemma}[Nominal weights and weights]\label{l:weights}
Let $D$ be admissible data for parameters $P=(N_0,d)$, $\sigma\in\PlD{D}$ and $j$ a natural number. Then
\begin{enumerate}[label=(\alph*)]
\item $\aD{D}/\scxo{D}\le\nw{D}{\sigma}\le\aD{D}/\scx{D}$; in particular $0\le\nw{D}{\sigma}$;
\item $\nw{D}{\sigma}+\scx{D}^{42}\le2\,\scb{D}$;
\item $\nw{D}{\sigma}\le\wt{D}{j}{\sigma}-\hhD{D}$;
\item if $j\le2\,\Dm{D}$, then $\wt{D}{j}{\sigma}\le\nw{D}{\sigma}+\scx{D}^{42}$, hence
$\wt{D}{j}{\sigma}\le2\,\scb{D}$, and $\log\wt{D}{j}{\sigma}\le\log\nw{D}{\sigma}+1/\scx{D}$.
\end{enumerate}
\end{lemma}
\begin{proof}
(a) $\scx{D}\le\scxo{D}$. (b) By (a) and Lemma~\ref{l:round},
$\nw{D}{\sigma}\le(\scx{D}\scb{D}+\ND{D})/\scx{D}=\scb{D}+\ND{D}/\scx{D}$, and
$\ND{D}/\scx{D}+\scx{D}^{42}\le\scb{D}$ by Theorem~B: an upper exponent of the left side is $42\lx{D}+1$.
(c) $\bone{D}\ge\aD{D}/\scx{D}$ and $\btwo{D}\ge\aD{D}/\scxo{D}$ by Lemma~\ref{l:round}. (d) By the same lemma the
difference $\wt{D}{j}{\sigma}-\nw{D}{\sigma}$ is at most $\ND{D}+2\,\Dm{D}\,\hD{D}+\hhD{D}$, with upper exponent
$40.5\,\lx{D}+33$, below $42\lx{D}$. With (b) this gives the second inequality. For the third,
$\nw{D}{\sigma}\,e^{1/\scx{D}}\ge\nw{D}{\sigma}+\nw{D}{\sigma}/\scx{D}$, and
$\scx{D}^{42}\le\nw{D}{\sigma}/\scx{D}$ by Theorem~B, since $42\lx{D}<\lb{D}-\lxo{D}-\lx{D}$.
\end{proof}
\leanline{\lean{nwicc}\sep\lean{nwnonneg}\sep\lean{nwadd}\sep\lean{nwlew}\sep\lean{wlenw}\sep\lean{wletwob}\sep\lean{logw}}

\sect{Chosen places and neat levels}{s:chosen}
\statement{ana}{\textbf{Inputs.}\enspace Let $D$ be admissible data for parameters $P=(N_0,d)$, and $I$ an index set.
(J1), unit~141: a finite place $w$ of $\ED{D}$ \emph{is chosen} if it contains $\cj{\Dpig{D}}$; for $i\in\{0,1\}$,
$\Ch{i}{D}$ is the set of chosen places that contain $\Dg{i}{D}$, and every chosen place lies in one of the two.
A place containing $\Dg{i}{D}$ has norm $\Dg{i}{D}$, and $\Dpig{D}\,\cj{\Dpig{D}}$ is $\Dg{i}{D}$ times a number
prime to $\Dg{i}{D}$. (J2), unit~141: the \emph{level} $\lvl{D}$ is a prime number with
$\lvl{D}\equiv1\bmod\ND{D}\,d$, and $\piN{D}$ is an integer of $\ED{D}$ with $\piN{D}\,\cj{\piN{D}}=\lvl{D}$.
(J3), unit~141: $\omega_D$ attaches to every finite place $w$ of $\ED{D}$ a character $\omega_D(w)$ of its residue
field, of order exactly $\Mom{P}{\Dmu{D}}$ if $w$ is chosen, and trivial otherwise. (J4), unit~122: $\Gf{D}{I}$ is
the unitary group of the standard hermitian form with index set $I$ over the finite adeles of $\ED{D}$,
$\Kz{D}{I}$ its subgroup of matrices that are integral with their inverses, $\Gam{D}{I}$ the image of the unitary
group over $\ED{D}$; for an integer $a$ of $\ED{D}$ the principal level of $a$ consists of the elements of
$\Kz{D}{I}$ congruent to $1$ modulo $a$. The general linear group with index set $I$ over a finite ring with $n$
elements has at most $n^{(\#I)^2}$ elements. (J5), unit~50, a lemma of Minkowski and Serre: if finitely many
distinct finite places of a CM field with more than one infinite place have norms whose product exceeds
$(2\,\#I)^{h}$, $h$ the degree of the field, then for every subgroup $K$ of the group of elements of the integral
subgroup that reduce to $1$ at these places, and every element $x$ of the adelic group, the intersection of the
rational group with $xKx^{-1}$ is trivial. The catalogue of the write-up refers for (J5) to Serre~\cite{serre}
and to Minkowski, as a special case, with the place not confirmed.}

\begin{lemma}[The chosen places]\label{l:chosen}
Let $D$ be admissible data for parameters $P=(N_0,d)$, $i\in\{0,1\}$, and $w$ a finite place of $\ED{D}$.
\begin{enumerate}[label=(\alph*)]
\item There is a bijection from $\PlD{D}$ onto $\Ch{i}{D}$ (the formal text fixes one); so $\Ch{i}{D}$ has $\fD{D}$
elements, and the product of the norms of all chosen places is $\gD{D}^{\fD{D}}$.
\item If $\Dg{i}{D}\in w$, then $\cj{w}$ is chosen if and only if $w$ is not. If $w$ is chosen, then $\cj{w}$ is
not chosen.
\item The family of all chosen places, indexed by the pairs of an index $i$ and an element of $\Ch{i}{D}$, has as
its range the set of chosen places; it is injective, and no member is the conjugate of a member.
\end{enumerate}
\end{lemma}
\begin{proof}
(a) The proposition of the first part on the places above a prime and an integer (the last of its section on Frobenius elements), for the prime $\Dg{i}{D}$, which is $\equiv1\bmod\ND{D}\,d$ by (D7), and the
integer $\cj{\Dpig{D}}$, with (J1); the places of $\Ch{i}{D}$ have norm $\Dg{i}{D}$. (b) The lemma of the first part on an integer whose norm is the prime times a number prime to it,
for the same prime and integer; a chosen place contains $\Dg{0}{D}$ or $\Dg{1}{D}$. (c) Places above
$\Dg{0}{D}\ne\Dg{1}{D}$ have different norms, and (b).
\end{proof}
\leanline{\lean{chequiv}\sep\lean{cardch}\sep\lean{prodq}\sep\lean{chstar}\sep\lean{notstar}\sep\lean{rangech}\sep
\lean{conjapart}}

\begin{lemma}[The level is unramified]\label{l:level}
Let $D$ be admissible data for parameters $P=(N_0,d)$ and $w$ a finite place of $\ED{D}$. If $\lvl{D}\in w$, then
$\lvl{D}$ has valuation exactly $1$ at $w$. If $\piN{D}\in w$, then $\piN{D}$ has valuation exactly $1$ at $w$.
\end{lemma}
\begin{proof}
The lemma of the first part on the absence of ramification and the lemma just used, for the prime $\lvl{D}$ of (J2) and the integer $\piN{D}$, whose norm is
$\lvl{D}$ times $1$.
\end{proof}
\leanline{\lean{valN}\sep\lean{valpi}}

\begin{definition}[The order and the level of the chosen places]\label{d:chosen}
Let $D$ be admissible data for parameters $P=(N_0,d)$. $\ordch{D}$ is the order of the product, over all chosen
places, of the general linear groups of rank $\Dr{D}\,\tD{D}$ over the residue fields. For an index set~$I$, which
is finite, $\Kch{D}{I}$ is the subgroup of the elements of $\Kz{D}{I}$ that reduce to $1$ at every chosen place.
\end{definition}
\leanline{\lean{ordch}\sep\lean{chlevel}}

\begin{lemma}\label{l:GL}
Let $D$ be admissible data for parameters $P=(N_0,d)$ and $I$ a finite index set~$I$. If $I$ has at most
$\Dr{D}\,\tD{D}$ elements, then the order of the product, over all chosen places, of the general linear groups with
index set $I$ over the residue fields divides $\ordch{D}$. If $a$ is an integer of $\ED{D}$ divisible by
$\cj{\Dpig{D}}$, then the principal level of $a$ is contained in $\Kch{D}{I}$.
\end{lemma}
\begin{proof}
Embed $I$ into a set with $\Dr{D}\,\tD{D}$ elements and extend matrices by the identity: this is an injective
homomorphism of groups, and the order of a subgroup divides the order of the group. Every chosen place contains
$\cj{\Dpig{D}}$, hence $a$.
\end{proof}
\leanline{\lean{cardGL}\sep\lean{princ}}

\begin{proposition}[Neat levels, and Minkowski's bound]\label{p:neat}
Let $D$ be admissible data for parameters $P=(N_0,d)$, $I$ a finite index set~$I$ with at most
$2\,\nD{D}{\Dm{D}}$ elements, and $K$ a subgroup~$K$ of $\Gf{D}{I}$.
\begin{enumerate}[label=(\alph*)]
\item If $K$ is contained in $\Kch{D}{I}$, then $\Gam{D}{I}\cap xKx^{-1}$ is trivial for every element~$x$ of
$\Gf{D}{I}$.
\item If $K$ is contained in $\Kz{D}{I}$, then the order of $\Gam{D}{I}\cap xKx^{-1}$ divides $\ordch{D}$ for
every element~$x$ of $\Gf{D}{I}$.
\end{enumerate}
\end{proposition}
\begin{proof}
(a) is (J5) for the family of chosen places, which is injective (Lemma~\ref{l:chosen}). Its hypothesis: the degree
of $\ED{D}$ is $2\fD{D}$ and the product of the norms is $\gD{D}^{\fD{D}}$, so it suffices that
$(4\,\nD{D}{\Dm{D}})^2<\gD{D}$, which holds by Theorem~B and Table~1: $11\lx{D}+32<20\lb{D}$. (b) Conjugation by
$x^{-1}$ followed by reduction at the chosen places is a homomorphism from $\Gam{D}{I}\cap xKx^{-1}$ to the product
of the general linear groups over the residue fields. Its kernel lies in $\Gam{D}{I}\cap x\,\Kch{D}{I}\,x^{-1}$,
which is trivial by (a). So the order divides that of the product, which divides $\ordch{D}$ by Lemma~\ref{l:GL},
since $2\,\nD{D}{\Dm{D}}\le\Dr{D}\,\tD{D}$ by Theorem~B.
\end{proof}
\leanline{\lean{neat}\sep\lean{mink}}

\begin{lemma}[Powers of the residue characters]\label{l:omega}
Let $D$ be admissible data for parameters $P=(N_0,d)$, $\nu$ an integer~$\nu$ with $0<|\nu|\le8\,\Dm{D}$, and $w$ a
finite place of $\ED{D}$. Then $\omega_D(w)^{\nu}\ne1$ if and only if $w$ is chosen.
\end{lemma}
\begin{proof}
At a place that is not chosen the character is trivial. At a chosen place its order $\Mom{P}{\Dmu{D}}$ is a
multiple of $\ND{D}\,d$ (the lemma of the first part on the arithmetic of the parameters, item (d)), and $8\,\Dm{D}<\ND{D}\,d$ by Theorem~B, so the order does
not divide $\nu$.
\end{proof}
\leanline{\lean{omega}. Two instances without names follow in the formal text: the last two lines of Table~1.}

For the line of $\log\ordch{D}$ in Table~1: by (J4) and Lemma~\ref{l:chosen}(a),
\[\ordch{D}\le\gD{D}^{\,\fD{D}\,(\Dr{D}\,\tD{D})^2},\]
and the exponent times $\log\gD{D}$ has an upper exponent below $32\lx{D}$.

\sect{Existence of admissible data}{s:exist}
\begin{proposition}[A construction of Behrend type]\label{p:behrend}
Let $\Gamma_1$, $\Gamma_2$ be commutative groups, $\psi\colon\Gamma_1\to\Gamma_2$ a homomorphism,
$\sigma_0,\gamma\in\Gamma_1$, and $n$, $M$, $L$ natural numbers with $L>0$ and $(LM)^n\le o$, where $o$ is the order
of $\psi(\gamma)$ (it is $0$ in the formal text if that order is infinite). Then there is a finite subset
$T\subseteq\Gamma_1$ such that
\begin{enumerate}[label=(\roman*)]
\item every element of $T$ is $\sigma_0\gamma^{\,i}$ for a natural number $i$;
\item $M^n/(n\,M^2)\le\#T$;
\item for all natural numbers $i$, $j$ with $0<i<j\le L$ and all $\sigma,\tau,\rho\in T$: if
$\psi(\tau/\sigma)^j=\psi(\rho/\sigma)^i$, then $\psi(\tau/\sigma)=1$.
\end{enumerate}
\end{proposition}
\begin{proof}
Among the $M^n$ vectors $u=(u_1,\dots,u_n)$ with integer coordinates $0\le u_k<M$, the sum of squares takes fewer
than $nM^2$ values; so for some value $q$ the sphere $S$ of the vectors with $\sum_ku_k^2=q$ has at least
$M^n/(nM^2)$ elements (this count is in the library of the prover). For a vector $u$ with natural coordinates put
$[u]=\sum_ku_k(LM)^{k-1}$; the map $u\mapsto[u]$ is additive, and on the vectors with all coordinates below $LM$ it
is injective with values below $(LM)^n\le o$. Hence $u\mapsto\psi(\gamma)^{[u]}$ is injective on these vectors.
Let $T$ be the set of the $\sigma_0\gamma^{[u]}$, $u\in S$. Then (i) holds, and (ii) because $L\ge1$ makes the
coordinates of $S$ smaller than $LM$. For (iii) let $\sigma,\tau,\rho$ come from $u,v,z\in S$ and write
$j=i+i'$ with $i'>0$. The hypothesis says $\psi(\gamma)^{(i+i')[v]}=\psi(\gamma)^{i[z]+i'[u]}$, and both exponents
are values of the bracket at vectors with coordinates at most $(i+i')(M-1)<LM$. So $(i+i')\,v=i\,z+i'\,u$: the
point $v$ lies between $z$ and $u$ on a segment, and all three lie on one sphere. Expanding
$i^2z_k^2=((i+i')v_k-i'u_k)^2$ and summing,
\[i\,(i\,q+i'q)=\sum_k i\,(i\,z_k^2+i'u_k^2)=\sum_k(i+i')\bigl(i'(u_k-v_k)^2+i\,v_k^2\bigr)
=(i+i')\,i'\sum_k(u_k-v_k)^2+(i+i')\,i\,q ,\]
so $u=v$, hence $\tau=\sigma$.
\end{proof}
\leanline{\lean{behrend}. The header of the Lean file cites Shapira~\cite{shapira}, Theorem~1, after
Behrend~\cite{behrend}; the artifact records these citations as found, with no place compared.}

\begin{lemma}[Pigeonhole in the class group]\label{l:pigeon}
Let $F$ be a number field~$F$ with class number $n$, and let $\mathfrak{p}_i$, $\mathfrak{q}_i$, for
$i\in\{0,\dots,n\}$, be finite places of $F$ such that $\mathfrak{p}_i\mathfrak{q}_i$ is principal for every $i$.
Then there are indices $i\ne j$ such that $\mathfrak{p}_i\mathfrak{q}_j$ is principal.
\end{lemma}
\begin{proof}
Two of the $n+1$ classes of the $\mathfrak{p}_i$ are equal, say those of $\mathfrak{p}_i$ and $\mathfrak{p}_j$ with
$i\ne j$; then the class of $\mathfrak{p}_i\mathfrak{q}_j$ is that of $\mathfrak{p}_j\mathfrak{q}_j$, the trivial
one.
\end{proof}
\leanline{\lean{pigeon}}

\statement{ana}{\textbf{Inputs.}\enspace (J6), unit~126, a weak form of the Siegel--Walfisz theorem: let $M>0$ be a
multiple of $4$ and $A$ a natural number. For all sufficiently large real $W$, every natural number $k>1$ with
$M^k\le W^A$ and every real $y$ with $W\le y\le W^A$, there is a prime number $p\equiv1\bmod M^k$ with
$y\le\log p\le y+W$. The catalogue of the write-up gives Montgomery and Vaughan~\cite{mv}, Corollary~11.19, of which
this is a weaker form. (J7), unit~106: for a prime $d\equiv7\bmod8$ and a prime number $p$ whose Legendre symbol
modulo $d$ is $1$, the ideal $(p)$ of $\Q(\sqrt{-d})$ is a product of two prime ideals of norm $p$; and the absolute
norm of a principal ideal $(y)$ is $y\,\cj y$. (J8), unit~106: if $q$ is an integer with $|q|>1$ and
$q^{o}\equiv1\bmod2^{\mu+1}$ for a natural number $o>0$, then $2^{\mu}\le(|q|+1)\,o$.}

\begin{proof}[Proof of Theorem A]
Write $Y=\log y$ and $W=y^{\varepsilon/6}=e^{\varepsilon Y/6}$, and let $\clno{P}$ be the class number of $\KP{P}$.
Fix a natural number $n$ with $n>6$ and $n\,\varepsilon/6>99\,N_0d$. Since $W$ exceeds every fixed multiple of every
fixed power of $Y$ when $y$ is large, we may assume all of the following: $Y\ge2$;
\[W\ge10^{12}\,Y^2,\qquad W\ge(8019N_0+2)\,10368^{\,n}\,Y^{n+1},\qquad W\ge8019\cdot2^n\,n\,Y,\]
\[W\ge2\,\clno{P}+10^6,\qquad W\ge\exp\bigl(10^{12}\,\CP{P}\bigr);\]
and the conclusion of (J6) for $M=N_0d$ and $A=n$ at this $W$.

\emph{Ranks and scales.} Choose natural numbers $i_2$, $i_3$ with $2^{i_2}\le e^{Y/2}<2^{i_2+1}$ and
$3^{i_3}\le e^{5Y}<3^{i_3+1}$, and put
\[m=2^{i_2+1},\qquad r=3^{i_3+1},\qquad\mu=22\,(i_2+1),\]
and take $Y$, $W$, $W^6$ for the logarithms of the three scales. Then (D3) and (D4) hold, and $2^\mu=m^{22}$. As
$i_2\log2\le Y/2$ and $Y\ge2$, we get $\mu\le16\,Y+22\le32\,Y$, which is (D2). For (D1):
$\log\log$ of the third scale is $\log W^6=\varepsilon Y$, which is at least $10^{12}\,\CP{P}$ by the last
assumption and at most $Y/100$; $10^6Y\le W$ and $10^6W\le W^6$ hold as $W\ge10^{12}Y^2$ and $W\ge10^6$; and
$6\log W=\varepsilon Y$. The same computations give the first four assertions of the theorem.

\emph{Primes.} First $(N_0d)^{2(\mu+6)}\le W^n$: the logarithm of the left side is at most
$2(32Y+6)\,N_0d\le99\,N_0d\,Y<n\log W$. Hence also $2^\mu\le W^n$. By (J6) with $k=2(\mu+6)$: for every real $y'$
with $W\le y'\le W^n$ there is a prime $p\equiv1$ modulo the square of $\Mom{P}{\mu}$ with
$y'\le\log p\le y'+W$. This gives (D9), because $W\le100\,W^6$, $2^\mu\le W^n$ and $W\le\sqrt{W^6}$. For
$i\in\{0,\dots,\clno{P}\}$ apply it to $y'=10W^6+2iW$, which is at most $W^7\le W^n$: there are primes $p_i\equiv1$
modulo the square of $\Mom{P}{\mu}$ with $10W^6+2iW\le\log p_i\le10W^6+(2i+1)W$. They increase strictly with $i$, and
$10W^6\le\log p_i\le11W^6$ because $(2\,\clno{P}+1)W\le W^2$. As $\Nn{P}{\mu}\,d$ divides $\Mom{P}{\mu}$, each $p_i$
is $\equiv1\bmod\Nn{P}{\mu}\,d$; so its Legendre symbol modulo $d$ is $1$ and $(p_i)=\mathfrak{p}_i\mathfrak{q}_i$ in
$\KP{P}$ with prime ideals of norm $p_i$ by (J7). By Lemma~\ref{l:pigeon} there are $i\ne j$ and an integer $\pi$
of $\KP{P}$ with $\mathfrak{p}_i\mathfrak{q}_j=(\pi)$, so that $\pi\,\cj\pi=p_ip_j$. With the two primes $p_i$, $p_j$
and the integer $\pi$, (D7) and (D8) hold.

\emph{Thick places.} Put $M=N_0\cdot3^{\lfloor\log_3\mu\rfloor+5}$, so that $\Nn{P}{\mu}=M\,2^\mu$, and let $\psi$ be
the class modulo $2^{\mu+1}$ on $\PlP{P}{\mu}$. The integers $-1$ and $M+1$ are prime to $N_0$; put
$\sigma_0=\pin{P}{\mu}{-1}$ and $\gamma=\pin{P}{\mu}{M+1}$. By the dictionary of the first part, item (d), $\psi(\gamma)$ is the
class of $M+1$, so its order $o$ satisfies $2^\mu\le(M+2)\,o$ by (J8). Choose a natural number $i_0$ with
$2^{n\,i_0}\le2^\mu/W<2^{n(i_0+1)}$; this is possible because $W\le e^{Y/2}<m\le2^\mu$. Now
$M\le243\,N_0(\mu+1)\le8019\,N_0Y$ and $324\mu\le10368\,Y$, so
\[(M+2)\,\bigl(324\mu\cdot2^{i_0}\bigr)^n\le(8019N_0+2)\,10368^{\,n}\,Y^{n+1}\cdot\frac{2^\mu}{W}\le2^\mu\le(M+2)\,o .\]
Thus Proposition~\ref{p:behrend} applies with $L=324\mu$ and the side $2^{i_0}$ in the role of its $M$. It gives a
set $T$ of places with (D6), with at least $2^{n\,i_0}/(n\,2^{2i_0})$ elements, all of the form
$\pin{P}{\mu}{-1}\cdot\pin{P}{\mu}{M+1}^{\,i}$. For a divisor $m'$ of $M$, such a place has class $-(M+1)^i=-1$
modulo $m'$, again by item (d) of the dictionary: this is (D5). So the record with thick set $T$ is admissible.

\emph{Many thick places.} It remains that $\twd{\mu}/W^6\le2^{n\,i_0}/(n\,2^{2i_0})$. Here
$\twd{\mu}\le243(\mu+1)2^\mu\le8019\,Y\cdot2^{n(i_0+1)}\,W$, and $2^{i_0}<W$ because $2^{n\,i_0}\le2^\mu/W<W^n$. So
\[\frac{\twd{\mu}}{W^6}\cdot\frac{n\,2^{2i_0}}{2^{n\,i_0}}\le\frac{8019\,Y\,2^{n}\,n\,W^{3}}{W^{6}}\le\frac{W^4}{W^6}\le1.\qedhere\]
\end{proof}
\leanline{\lean{exdata}}

\sect{Notes on the formal text}{s:notes}
\subhead{Names and arguments.} In the formal text the scales, the ranks $\kD{D}$, $\tD{D}$, the field and every
other object of Definitions~\ref{d:scales} to~\ref{d:sums} are functions of the data, and the note writes them so.
The components of the record are written in the same way, $\Dm{D}$, $\Dr{D}$, $\Dmu{D}$, $\thick{D}$,
$\Dg{i}{D}$, $\Dpig{D}$; the overview of the write-up writes some of them without the argument. The record stores
the logarithms of the scales, and the scales are their exponentials: the note writes $\lx{D}$ for the stored
number. The component that the formal text calls \texttt{s} is $\Dmu{D}$ here. In Lean the hypotheses
$\hlu{D}(F;\,l)$ and $\hll{D}(G;\,l')$ of Theorem~B and of Lemmas~\ref{l:rules} to~\ref{l:round} are instance
arguments: Lean finds them by search, from the rules of Proposition~\ref{p:table} and the lines of Table~1, each of
which is an instance; and the hypothesis $l<l'$ of Theorem~B has a default proof by a decision procedure.

\subhead{What was not redone by hand.} The formal text proves each line of Table~1, and each use of Theorem~B, by
one call of a decision procedure on the conditions (D1). In this note the inequalities between exponents are
written out for the lines of $\Dm{D}$, $\twd{\Dmu{D}}$, $\fD{D}$ and for the uses in Lemmas~\ref{l:ranks},
\ref{l:weights}, \ref{l:omega} and Proposition~\ref{p:neat}; the other lines of the table are stated as the formal
text has them. The exponents given in the proofs are mine and may differ from those that Lean finds. The proof
of Theorem~A is complete here; the formal proof closes several of its numerical steps by automation, and which
goal each tactic settles there was inferred by the reader of the unit from the order of the conditions, not from
Lean's output.

\subhead{Signatures that could not be settled without Lean.} (1) For five declarations treated here the list of
arguments is derived by the reader of the unit and not printed by Lean: Lemma~\ref{l:cge}, the range and the
conjugation statement of Lemma~\ref{l:chosen}(c), and the lines of $\fD{D}$ and of $\CP{P}$ in Table~1. I read them
in the Lean text and agree: the letter that looks like a section variable is in each case another object (the
constant of the parameters, written with a dot after $P$; a variable bound inside the statement). (2) The two
lines of Table~1 for $\wt{D}{0}{\sigma}$ and $\wt{D}{j}{\sigma}$ are instances that stand, in the Lean file, after a
command that includes the hypothesis $j\le2\,\Dm{D}$ in what follows. Lean's documentation says that this command
concerns theorems and no other declarations; and the proof of Lemma~\ref{l:weights}(d) needs the second instance
to be found by search, which could not supply such a hypothesis. So I read both instances without the hypothesis,
and the first one without $j$, against the census of the workspace; only Lean can settle it. (3) The rule for
quotients is an instance that takes bounds for $F$ and for $1/G$. The census lists a further bound for $G$ as a
possible argument; by Lean's rule for definitions, which keeps a section variable only if the term uses it, it is
not one. This too is a reading of the rule, not Lean's answer.

\subhead{Where the formal statements differ from the natural ones.} (D5) quantifies over a number $m$ and a proof
that it divides the order of the root of unity; the name $m$ hides the component $\Dm{D}$ there. The subtraction in
$\kD{D}$, $\posD{D}{j}$ and $\rhg{D}$ is that of natural numbers; Lemma~\ref{l:ranks} shows that it is exact for
$\kD{D}$. In (D7) the formal text says that the pair of primes is injective and that the square of
$\Mom{P}{\Dmu{D}}$ divides $\Dg{i}{D}-1$ in natural numbers. In Theorem~A the hypothesis is $0<\varepsilon/6$, and
``for every sufficiently large $y$'' is a statement along the filter at infinity. In
Proposition~\ref{p:behrend} the bound (ii) is a real quotient, which Lean reads as $0$ when $n=0$ or $M=0$, and the
order of an element of infinite order is $0$; then the hypothesis forces $M=0$ and the statement is empty.
Lemma~\ref{l:rules}(a) is all that the formal text uses of positivity. The definition of $\rhg{D}$ and of
$\posD{D}{j}$ are not used in this unit.

\subhead{How this note was written.} From the dossier of the unit: the texts of its members and the structured
proofs of the five files treated here, of which the longest is the proof of Theorem~A. I read the five Lean files
completely and compared every statement of this note with the Lean statement, argument by argument. The proof of
Theorem~A follows the formal one step by step, with its constants; the proofs of Lemmas~\ref{l:ranks}
and~\ref{l:weights} and of Proposition~\ref{p:table} replace calls of automation by explicit exponents. Of the
Inputs I opened only the comments of the formal text and the reading of the dossier, not the proofs. The references
are those of the catalogue of the write-up and of the artifact's own list; I did not open the cited works. A second
reading of the unit by the reading agent arrived while this note was written; I compared its remarks on this part
with the Lean text, and they agree with what is stated here. The
program check of the unit's dossier covered \UnitSteps{} leaf steps of \UnitProofs{} proofs.

\head{Appendix A. Concordance with the Lean artifact}
The formal names of the objects: \texttt{D.logX}, \texttt{D.logX₁}, \texttt{D.logB} are $\lx{D}$, $\lxo{D}$,
$\lb{D}$; \texttt{D.s}, \texttt{D.m}, \texttt{D.r} are $\Dmu{D}$, $\Dm{D}$, $\Dr{D}$; \texttt{D.g i} is $\Dg{i}{D}$ and
\texttt{D.piG} is $\Dpig{D}$; \texttt{D.Φ} is $\PlD{D}$; \texttt{P.c} is $\CP{P}$; \texttt{primeSum} is the prime
sum, \texttt{dirCost} the cost of one direction. Three statements as the artifact has them:
\quoted{Data}
\quoted{exdata}
\quoted{leBd}
{\fontsize{8.6}{11.6}\selectfont\renewcommand{\leanfont}{\ttfamily\fontsize{7.7}{10}\selectfont\color{leancol}}
\setlength{\LTleft}{0pt}\setlength{\LTright}{\fill}\setlength{\tabcolsep}{0pt}
\begin{longtable}{@{}p{27mm}p{62mm}p{40mm}r@{}}
\textbf{In this note} & \textbf{Lean declaration} & \textbf{File} & \textbf{Lines}\\\midrule\endhead
\ConcordanceRows
\end{longtable}}

\head{Notation}
\begingroup\fontsize{8.6}{10.9}\selectfont\raggedright
\setlength{\columnsep}{9mm}\setlength{\parskip}{0pt}
\begin{multicols}{2}
\nt{$P=(N_0,d)$}{parameters}
\nt{$D$}{admissible data for $P$ (Definition~\ref{d:data})}
\nt{$\scx{D}$, $\scxo{D}$, $\scb{D}$}{the three scales}
\nt{$\Dmu{D}$, $\Dm{D}$, $\Dr{D}$}{the exponent, the number of lifts, the rank of the rows}
\nt{$\thick{D}$}{the set of thick places}
\nt{$\Dg{i}{D}$, $\gD{D}$, $\Dpig{D}$}{the two auxiliary primes, their product, an integer of that norm}
\nt{$\CP{P}$, $C_{\mathrm{jet}}$}{the constant of the parameters and one of its summands}
\nt{$\twd{\mu}$, $\Nn{P}{\mu}$, $\Mom{P}{\mu}$}{as in the first part}
\nt{$\KP{P}$, $\chiK{P}$, $\EP{P}{\mu}$, $\PlP{P}{\mu}$}{as in the first part}
\nt{$\pin{P}{\mu}{\nu}$, $\pcl{P}{m\mid\Nn{P}{\mu}}$}{the place named by $\nu$; the class modulo $m$}
\nt{$\OmP{P}{\mu}$}{the period}
\nt{$\ND{D}$, $\ED{D}$, $\PlD{D}$}{order, field and places for the exponent of $D$}
\nt{$\fD{D}$, $\rdD{D}$}{the number of infinite places; the root discriminant}
\nt{$\pinD{D}{j}$, $\tops{D}$}{the place named by $j$; the tops}
\nt{$\kD{D}$, $\tD{D}$, $\nD{D}{j}$, $\posD{D}{j}$}{ranks}
\nt{$\hhD{D}$, $\hD{D}$, $\rhg{D}$}{the half-gap, the gap, the reduced half-gap}
\nt{$\rup{D}{y}$}{$y$ rounded up to a multiple of $\ND{D}$}
\nt{$\aD{D}$, $\bone{D}$, $\btwo{D}$}{the order and the two rounded quotients}
\nt{$\nw{D}{\sigma}$, $\wt{D}{j}{\sigma}$}{nominal weight; weight}
\nt{$\Ask{D}{y}{\sigma}$, $\lambda_{d,1}$, $\Lambda$}{the prime sum; its coefficients; von Mangoldt's function}
\nt{$\budg{D}$, $\jet{D}{j}$, $\UD{D}$}{budget; jet size; cost of one direction}
\nt{$\hlb{D}{\ast}{F}{l}$, $\hlu{D}$, $\hll{D}$}{bounds on the logarithmic scale; $\ast\in\{\top,\bot\}$}
\nt{$\lvl{D}$, $\piN{D}$}{the level and an integer of that norm}
\nt{$\Ch{i}{D}$, $\omega_D$}{chosen places above $\Dg{i}{D}$; residue characters}
\nt{$\Gf{D}{I}$, $\Kz{D}{I}$, $\Gam{D}{I}$, $\Kch{D}{I}$}{adelic unitary group, integral subgroup, rational points, level of the chosen places}
\nt{$\ordch{D}$}{the order of Definition~\ref{d:chosen}}
\nt{$\clno{P}$}{the class number of $\KP{P}$}
\nt{$F,G,H$}{real numbers; $F$ also a number field in Lemma~\ref{l:pigeon}}
\nt{$l,l',l''$}{exponents}
\nt{$A,B$}{natural numbers; $A$ also the exponent in (J6)}
\nt{$y,y'$}{real numbers; $y$ the scale in Theorem~A}
\nt{$Y,W$}{$\log y$ and $y^{\varepsilon/6}$ in the proof of Theorem~A}
\nt{$\varepsilon$}{the exponent of Theorem~A}
\nt{$\sigma,\tau,\rho$}{infinite places, or elements of the set $T$}
\nt{$\nu$}{an integer}
\nt{$i,j,k,n,o,q$}{natural numbers and integers: indices, exponents, orders}
\nt{$i',i_0,i_2,i_3$}{exponents in two proofs}
\nt{$m,r,\mu,m'$}{in the proof of Theorem~A, the components under construction and a divisor; $m$ also a divisor in (D5)}
\nt{$p,p_i$}{prime numbers}
\nt{$M,L$}{a side and a bound (Proposition~\ref{p:behrend}); $M$ also a modulus}
\nt{$\Gamma_1,\Gamma_2,\psi,\sigma_0,\gamma$}{groups, a homomorphism, two elements; $\psi$ also a class map}
\nt{$T,S$}{a finite set of group elements; a sphere of integer vectors}
\nt{$u,v,z$, $[u]$}{integer vectors and the bracket of the proof of Proposition~\ref{p:behrend}}
\nt{$\mathfrak{p}_i,\mathfrak{q}_i$, $\pi$}{prime ideals; a generator}
\nt{$w$, $\cj{w}$, $\Nm(w)$}{a finite place, its conjugate, its norm}
\nt{$I$, $K$, $x$, $a$, $h$}{an index set, a subgroup, a group element (also a root of unity in Definition~\ref{d:C}), an integer, a degree}
\nt{$\varphi$, $\pi$, $e$, $L$}{Euler's function; the circle constant in $2\pi$; the base of $\exp$; in ``$L$-function''}
\nt{$\mathrm{cs}(d)$, $\mathrm{rd}$, $\Cy{N_0}$}{as in the first part}
\end{multicols}
\endgroup

\head{References}
\begingroup\fontsize{9.4}{12.6}\selectfont
\begin{thebibliography}{9}
\bibitem{behrend} F. A. Behrend, Proc. Nat. Acad. Sci. 32 (1946). As cited by the artifact; not opened.
\bibitem{mv} H. L. Montgomery, R. C. Vaughan, \emph{Multiplicative Number Theory I. Classical Theory}, Cambridge
University Press, 2007. Corollary~11.19, as recorded in the catalogue of the write-up.
\bibitem{serre} J.-P. Serre, Rigidité du foncteur de Jacobi d'échelon $n\ge3$, Séminaire Henri Cartan 13 (1960/61),
appendix. As recorded in the catalogue of the write-up, place not confirmed.
\bibitem{shapira} A. Shapira, Behrend-type constructions for sets of linear equations, Acta Arith. (2006),
Theorem~1. As cited by the artifact; not opened.
\end{thebibliography}
\endgroup

\end{document}
